% !TEX root = ../main.tex
\chapter{Regression Analysis}\label{ch:regression}
\begin{paracol}{2}

\begin{sources}
\begin{tabularx}{\linewidth}{@{}l l L@{}}
\toprule
\textbf{Lecture} & \textbf{Speaker} & \textbf{Content}\\
\midrule
2024 L6 (from 1:07:22) & P.~Kempthorne & Multiple linear regression: set-up, polynomial, Fourier and autoregressive fits of the high-yield spread, the steps of model fitting, assumptions on the errors.\\
2024 L8 & P.~Kempthorne & OLS, hat matrix and geometry; normal linear model via moment-generating functions; $t$ and $F$ tests; prostate-cancer example (standardised covariates, $R^2$, diagnostics); Gauss--Markov; GLS.\\
2024 L11 & P.~Kempthorne & Recap, $t^2=F$, GLS, maximum likelihood, M-estimators (robust, quantile), ridge, principal-components regression, lasso; ETF case study; CAPM regressions (GE, all S\&P~500 stocks) and tests for a change of regime.\\
2013 L6 & P.~Kempthorne & The same theory: OLS, Gauss--Markov with proof, GLS, distribution theory via MGFs and the QR decomposition, maximum likelihood.\\
\bottomrule
\end{tabularx}

\smallskip
\textbf{Files.} Slides \file{mit18_642_f24_lec06_1.pdf} (used in L6, L8 and L11);
\file{mit18_642_f24_lec06_2.pdf} and \file{mit18_642_f24_lec06_3.pdf} (polynomial and Fourier fits
of the high-yield spread); \file{mit18_642_f24_lec06_4.pdf}, the knitted \file{ETF_Casestudy.rmd}
of \file{mit18_642_f24_etf_casestudy.zip} (which also contains \file{ETF_Casestudy_XLB.rmd/.pdf},
\file{spyders_symbols_labels.csv}, \file{casestudy_1_0_etfs.RData}); note \emph{Hypothesis Testing in
Regression Analysis} \file{mit18_642_f24_lec08.pdf}; 2013 notes \file{MIT18_S096F13_lecnote6.pdf};
2013 Case Study~1 \file{CaseStudy1.pdf} with \file{fm_casestudy_1_0.r} and
\file{fm_casestudy_1_rcode.r}; 2013 Case Study~2 \file{CaseStudy2.pdf} with
\file{fm_casestudy_2_rcode.r}, \file{fm_casestudy_fx_1.r} and \file{fred_fxrates.txt}. Installation
scripts and R instructions are described in \cref{app:rsetup}.
\textbf{Problem sets.} 2013 PS3 (all problems). Least squares via the SVD (2013 PS1, B-4) is in
\cref{ch:linalg}.
\textbf{Not published.} The note on CAPM regressions for all S\&P~500 stocks shown in 2024 L11; only
its GE part is in \file{lec08} (see \cref{ch05:capm}).
\end{sources}

\section*{Motivation}
Linear regression is the most used statistical tool in finance: every beta, hedge ratio, factor
exposure, CAPM test or ``is this strategy's alpha real?'' question is a regression. The mathematics
has two layers, and this chapter develops both:
\begin{itemize}
  \item \emph{Geometry} (no probability): least squares is the orthogonal projection of the data
        vector $\bm y\in\R^n$ on the column space of the design matrix $X$ (\cref{ch05:sec-ols}).
        This is linear algebra from \cref{ch:linalg}.
  \item \emph{Statistics}: under assumptions on the errors, the projection has optimality
        properties (Gauss--Markov, \cref{ch05:sec-gm}) and known sampling distributions
        (\cref{ch05:sec-normal}), which give $t$ and $F$ tests (\cref{ch05:sec-tests}), including
        tests for a change of regime (\cref{ch05:sec-chow}).
\end{itemize}
The rest is about what to do when the assumptions fail: diagnostics (\cref{ch05:sec-diag}), robust
and quantile estimators (\cref{ch05:sec-mle}), and shrinkage (\cref{ch05:sec-shrink}). The course's
message, repeated in every lecture, is methodological: propose a model, fit it, \emph{check the
assumptions}, revise \ts{24}{6}{1:18:03}. The case studies (sector ETFs, CAPM for GE and the S\&P~500,
the Chinese yuan) show what this looks like on real data.

% ------------------------------------------------------------------
\section{The linear regression model}\label{ch05:sec-model}

\paragraph{Data and goals.} We observe $n$ cases $i=1,\dots,n$; for each, a \emph{response}
(dependent variable) $y_i$ and $p$ \emph{explanatory} (independent) variables
$\bm x_i=(x_{i,1},\dots,x_{i,p})\T$ \ts{24}{6}{1:07:22}. In finance $y_i$ may be a stock return and
$\bm x_i$ market or macroeconomic factors \ts{13}{6}{2:04}. The goals are prediction, causal inference
(which factors really drive $y$?), approximation of an unknown function, and the discovery or validation
of functional relationships \ts{24}{6}{1:08:24}.

\begin{definition}[Linear model]\label{ch05:def-linmodel}
The \emph{general linear model} specifies the conditional distribution of $y_i$ given $\bm x_i$ as
\[
  y_i = \hat y_i + \varepsilon_i,\qquad \hat y_i=\beta_1x_{i,1}+\beta_2x_{i,2}+\dots+\beta_px_{i,p},
\]
with regression parameters $\bm\beta=(\beta_1,\dots,\beta_p)\T$ constant over the cases and a residual
(error) $\varepsilon_i$ varying over the cases. In matrix form
\[
  \bm y = X\bm\beta+\bm\varepsilon,\qquad
  \bm y=\begin{pmatrix}y_1\\ \vdots\\ y_n\end{pmatrix},\quad
  X=\begin{pmatrix}x_{1,1}&\cdots&x_{1,p}\\ \vdots&\ddots&\vdots\\ x_{n,1}&\cdots&x_{n,p}\end{pmatrix}\in\R^{n\times p}:
\]
cases are rows, variables are columns \ts{13}{6}{29:24}. An intercept is included by taking a first
column of ones.
\end{definition}

\begin{coursenote}[Erratum in the slides]
Slides 3 and 7 of \file{lec06_1} (and the 2013 notes) write the last term as $\beta_{i,p}x_{i,p}$: it must
be $\beta_px_{i,p}$, since the parameters do not depend on the case; Kempthorne corrects this in class
\ts{24}{6}{1:09:30}. In the same slides the bottom-right entry of $X$ is printed $x_{p,n}$ instead of
$x_{n,p}$.
\end{coursenote}

\paragraph{``Linear'' means linear in $\bm\beta$.} The columns of $X$ can be any functions of the
underlying variables \ts{24}{6}{1:10:32} \ts{13}{6}{6:19}:
\begin{itemize}
  \item polynomials, $x_{i,j}=x_i^{\,j}$ (a truncated Taylor expansion);
  \item Fourier terms, $x_{i,j}=\sin(jx_i)$ or $\cos(jx_i)$, for cyclical behaviour;
  \item time-series regressions, where $i$ indexes time and the regressors include lagged values
        $y_{i-1},y_{i-2},\dots$ of the response (autoregressions, \cref{ch:timeseries}) \ts{13}{6}{7:24}.
\end{itemize}
What matters is that $\hat y_i$ is linear in the parameters; this is what makes the estimation a problem
of linear algebra.

\begin{lab}[high-yield spread, functions of time against persistence]{ch05:lab-hyspread}
\textbf{Question.} Can the \emph{high-yield spread} be described by a function of calendar time, and what is
actually predictable in it? (Case study of \file{lec06_2} and \file{lec06_3}.)

\textbf{Data.} \file{data/fm_hyspread.dat}: daily spread 2010--2017, Moody's Baa corporate yield minus the 10-year US
Treasury yield, in percentage points \ts{24}{6}{1:11:33} (FRED series BAA10Y; 2001 days). Lower-rated borrowers pay
more than the Treasury, so it is always positive.

\textbf{In Python} (numpy, least squares with \texttt{np.linalg.lstsq}):
\begin{enumerate}[1.]
  \item Plot the spread. Regress it on powers of time, from the constant to an 8th-order polynomial, and record
        $R^2$ (rescale time to $[-1,1]$ first: why?).
  \item Regress it on Fourier terms $\sin(ju),\cos(ju)$, $j=1,\dots,J$, with $u$ running from 0 to $2\pi$ over the
        sample, for $J=1,\dots,4$.
  \item Regress today's spread on yesterday's, $y_t=a+by_{t-1}+\varepsilon_t$. Compare the three $R^2$, look at the
        fitted curves at the ends of the sample, and say which model you would use to forecast tomorrow.
\end{enumerate}
Reference script: \file{labs/ch05_hyspread.py} (\texttt{--plot} draws the fits).
\end{lab}
\begin{solution}
\textbf{Course results.} Higher orders follow the two big waves (peaks around 2012 and early 2016) better and
better, but the polynomials swing wildly at the ends of the sample, and nobody believes the spread
\emph{is} an 8th-order polynomial in calendar time \ts{24}{6}{1:12:37}. In class a third model was shown
(not in the files): the \emph{autoregression} of today's spread on yesterday's, whose fit is almost
indistinguishable from the data, with $R^2=0.995$ \ts{24}{6}{1:13:44}. (In \file{lec06_3} the legend colours do not
match the curves: the 1-cycle fit is drawn in blue, not green.)

\textbf{With the FRED data} (the course's exact window is not given, so the numbers differ a little):
\begin{center}\small
\begin{tabular}{@{}lcccccc@{}}
\toprule
model & poly 1 & poly 4 & poly 8 & Fourier 2 & Fourier 4 & AR(1)\\
\midrule
$R^2$ & 0.18 & 0.51 & 0.81 & 0.59 & 0.82 & 0.997\\
\bottomrule
\end{tabular}
\end{center}
Rescaling time avoids powers of numbers like 2017, whose columns are almost collinear (an ill-conditioned
$X\T X$). The AR(1) slope is $b=0.9997$: the spread is close to a random walk.

\textbf{Lesson.} A good in-sample fit is not a model. Deterministic functions of time only describe the
past; the autoregression exploits \emph{persistence} (tomorrow's spread is close to today's), which is
what is actually predictable. But with $b\approx1$ its $R^2$ says little: ``tomorrow = today'' already has
$R^2\approx0.997$; the honest comparison is on the daily \emph{changes} (\cref{ch:timeseries}).
\end{solution}

\paragraph{Steps for fitting a model} \ts{24}{6}{1:14:47} \ts{13}{6}{8:24}.
\begin{enumerate}[1.]
  \item Propose a model: the response and its \emph{scale} (e.g.\ log prices, so that errors are
        percentage moves \ts{13}{6}{9:25}), the explanatory variables (including functions of them), and
        assumptions on the distribution of $\bm\varepsilon$.
  \item Define a criterion for judging estimators (least squares, maximum likelihood, robust, Bayes,
        methods for missing data) \ts{13}{6}{23:07}.
  \item Characterise the best estimator under that criterion and compute it.
  \item Check the assumptions of step 1: residual analysis (the errors $\varepsilon_i$ are unobservable, the
        residuals $\hat\varepsilon_i=y_i-\hat y_i$ are their proxies), influence diagnostics, outlier detection
        \ts{13}{6}{26:16}.
  \item If necessary modify the model and/or the assumptions and go back to step 1.
\end{enumerate}
Step 5 is where the work is on hard problems: the model should be tailored to the process, not the
other way round \ts{13}{6}{12:31}.

\begin{definition}[Assumptions on the errors]\label{ch05:def-assumptions}
In increasing generality \ts{24}{6}{1:15:54} \ts{13}{6}{14:38}:
\begin{itemize}
  \item \emph{Normal linear model}: $\varepsilon_1,\dots,\varepsilon_n$ i.i.d.\ $N(0,\sigma^2)$.
  \item \emph{Gauss--Markov}: $\E[\bm\varepsilon]=\bm 0$ and $\Cov(\bm\varepsilon)=\sigma^2I_n$, i.e.\ zero mean,
        constant variance (homoscedastic), uncorrelated. Uncorrelated is weaker than independent
        \ts{13}{6}{15:40}.
  \item \emph{Generalised Gauss--Markov}: $\E[\bm\varepsilon]=\bm 0$, $\Cov(\bm\varepsilon)=\sigma^2\Sigma$ with a general
        positive definite $\Sigma$ (correlated and/or heteroscedastic errors).
  \item \emph{Non-Gaussian} errors: Laplace, Pareto, or \emph{contaminated normal}: a fraction $1-\delta$ of the
        $\varepsilon_i$ i.i.d.\ $N(0,\sigma^2)$ and a fraction $\delta$ from some contaminating distribution.
\end{itemize}
\end{definition}
Correlated residuals arise naturally when a straight line is fitted to a curved relationship
(residuals are systematically above, then below the line) or when unobserved factors persist in time
\ts{13}{6}{20:58}. Generalised Gauss--Markov is the typical situation of time-series regressions, and
heavy-tailed errors that of asset returns \ts{24}{6}{1:16:57}.

\begin{intuition}[A physicist's reading]
This is curve fitting with error bars. Least squares with equal error bars is $\chi^2$ minimisation;
Generalised Gauss--Markov is the case of a full covariance matrix of the measurements; the five steps are
the usual loop ``model, fit, look at the residuals (pulls), improve the model''. The one difference from a
lab: in finance there is no repeated experiment, and the ``error bars'' themselves must be estimated from
the same data.
\end{intuition}

% ------------------------------------------------------------------
\section{Ordinary least squares and its geometry}\label{ch05:sec-ols}

\subsection{Normal equations}
The \emph{least-squares criterion} is \ts{24}{8}{1:06}
\[
  Q(\bm\beta)=\sum_{i=1}^n\bigl(y_i-\bm x_i\T\bm\beta\bigr)^2=(\bm y-X\bm\beta)\T(\bm y-X\bm\beta)=\|\bm y-X\bm\beta\|^2,
\]
and the \emph{ordinary least-squares} (OLS) estimate $\hat{\bm\beta}$ minimises it. Differentiating,
\[
  \frac{\partial Q}{\partial\beta_j}=\sum_{i=1}^n2(-x_{i,j})\bigl(y_i-\bm x_i\T\bm\beta\bigr)=-2X_{[j]}\T(\bm y-X\bm\beta),
  \qquad
  \frac{\partial Q}{\partial\bm\beta}=-2X\T(\bm y-X\bm\beta),
\]
where $X_{[j]}$ is the $j$-th column of $X$ \ts{13}{6}{33:40}. The Hessian is
$\partial^2Q/\partial\bm\beta\,\partial\bm\beta\T=2X\T X$, positive semi-definite since
$\bm v\T X\T X\bm v=\|X\bm v\|^2\ge0$: $Q$ is convex, so every stationary point is a global minimum
\ts{24}{8}{5:16} \ts{13}{6}{32:34}. Setting the gradient to zero gives the \emph{normal equations}
\begin{equation}\label{ch05:eq-normal}
  X\T(\bm y-X\hat{\bm\beta})=\bm 0\iff X\T X\hat{\bm\beta}=X\T\bm y
  \quad\Longrightarrow\quad \hat{\bm\beta}=(X\T X)^{-1}X\T\bm y .
\end{equation}
$X\T X$ is invertible iff $X\bm v=\bm 0\Rightarrow\bm v=\bm 0$, i.e.\ iff $X$ has \emph{full column rank} $p$
\ts{24}{8}{7:33}: no explanatory variable may be a linear combination of the others \ts{13}{6}{37:50}.
Without full rank $\hat{\bm\beta}$ is not unique, but the fitted values are (see below), which is enough
for pure prediction \ts{13}{6}{37:50}.

\begin{remark}
Kempthorne jokes that there is nothing ``normal'' about the normal equations \ts{24}{8}{7:33}. A useful
way to remember the name: \eqref{ch05:eq-normal} says that the residual vector is \emph{normal}
(perpendicular) to every column of $X$.
\end{remark}

\subsection{The hat matrix}
The fitted values are $\hat{\bm y}=X\hat{\bm\beta}=H\bm y$ with the $n\times n$ \emph{hat matrix}
\ts{24}{8}{8:36} \ts{13}{6}{39:58}
\[
  H=X(X\T X)^{-1}X\T ,
\]
and the residuals are $\hat{\bm\varepsilon}=\bm y-\hat{\bm y}=(I_n-H)\bm y$.

\begin{proposition}[Properties of the hat matrix]\label{ch05:prop-hat}
Let $X\in\R^{n\times p}$ have full column rank.
\begin{enumerate}[(i)]
  \item $H\T=H$ and $H^2=H$; the same holds for $I_n-H$.
  \item $H$ is the orthogonal projection onto $\operatorname{col}(X)$, and $I_n-H$ the orthogonal projection
        onto $\operatorname{col}(X)^\perp$.
  \item $\tr H=\rank H=p$ and $\tr(I_n-H)=n-p$.
  \item $X\T\hat{\bm\varepsilon}=\bm 0$, hence $\hat{\bm\varepsilon}\perp\hat{\bm y}$ and
        $\|\bm y\|^2=\|\hat{\bm y}\|^2+\|\hat{\bm\varepsilon}\|^2$.
\end{enumerate}
\end{proposition}
\begin{proof}
(i) $H\T=X[(X\T X)^{-1}]\T X\T=H$ because $X\T X$ is symmetric;
\[H^2=X(X\T X)^{-1}(X\T X)(X\T X)^{-1}X\T=H\]
\ts{24}{8}{10:41}; $(I-H)^2=I-2H+H^2=I-H$ \ts{24}{8}{11:45}.

(ii) If $\bm v=X\bm a\in\operatorname{col}(X)$ then $H\bm v=X(X\T X)^{-1}X\T X\bm a=\bm v$; if $\bm w\perp\operatorname{col}(X)$ then
$X\T\bm w=\bm 0$ and $H\bm w=\bm 0$. Since $\R^n=\operatorname{col}(X)\oplus\operatorname{col}(X)^\perp$, $H$ is the orthogonal
projection, and $I-H$ the complementary one.

(iii) $\tr H=\tr\bigl((X\T X)^{-1}X\T X\bigr)=\tr I_p=p$, using $\tr(AB)=\tr(BA)$; for a projection, trace = rank.

(iv) is \eqref{ch05:eq-normal}; then $\hat{\bm y}\T\hat{\bm\varepsilon}=\hat{\bm\beta}\T X\T\hat{\bm\varepsilon}=0$ and Pythagoras
\ts{24}{8}{14:54}.
\end{proof}

\begin{figure}[ht]
\centering
\begin{tikzpicture}[scale=1.05,>=Stealth]
  \fill[defcol!8] (0,0) -- (6.4,0) -- (8.2,2.3) -- (1.8,2.3) -- cycle;
  \draw[defcol!60] (0,0) -- (6.4,0) -- (8.2,2.3) -- (1.8,2.3) -- cycle;
  \node[defcol,anchor=west] at (6.55,0.35) {$\operatorname{col}(X)$};
  \coordinate (O) at (1.2,0.6);
  \coordinate (Yh) at (5.0,1.4);
  \coordinate (Y) at (5.0,4.3);
  \draw[->,thick,black!60] (O) -- (3.4,0.6) node[below,black]{\small $X_{[1]}$};
  \draw[->,thick,black!60] (O) -- (2.3,1.85) node[left,black]{\small $X_{[2]}$};
  \draw[->,very thick,intcol] (O) -- (Y) node[above]{$\bm y$};
  \draw[->,very thick,defcol] (O) -- (Yh) node[below right]{$\hat{\bm y}=H\bm y$};
  \draw[->,very thick,thmcol] (Yh) -- (Y) node[midway,right]{$\hat{\bm\varepsilon}=(I-H)\bm y$};
  \draw[black!70] (5.0,1.7) -- (4.707,1.638) -- (4.707,1.338);
  \fill (O) circle (1.2pt) node[below left]{\small $\bm 0$};
\end{tikzpicture}
\caption{Least squares in $\R^n$: $\hat{\bm y}$ is the orthogonal projection of $\bm y$ on the column space of
$X$ (spanned here by two columns), and the residual is perpendicular to it
\ts{24}{8}{13:50} \ts{13}{6}{42:08}.}\label{ch05:fig-geometry}
\end{figure}

\begin{intuition}[Two pictures of the same fit]
In the usual scatter plot we look at $n$ points in $\R^{p}$ and a fitted hyperplane. In
\cref{ch05:fig-geometry} we look at \emph{one} point $\bm y$ in $\R^n$ (one axis per case) and a
$p$-dimensional subspace: least squares drops the perpendicular \ts{13}{6}{42:08}. The second picture is
the one in which all the distribution theory below becomes transparent: the residual lives in an
$(n-p)$-dimensional space, which is why ``$n-p$ degrees of freedom'' appear everywhere.
\end{intuition}

\begin{exercise}[PS3 (2013), Problem 1(a)--(c): the hat matrix]\label{ch05:ex-ps3-1ac}
\leavevmode
\begin{enumerate}[(a)]
  \item Prove that $H$ is a projection matrix: $H\T=H$ and $H^2=H$.
  \item The diagonal element $H_{ii}$ is the \emph{leverage} of case $i$. Show that $\partial\hat y_i/\partial y_i=H_{ii}$.
  \item Show that the average of the $H_{ii}$ is $p/n$. \emph{Hint:} $\tr(AB)=\tr(BA)$.
\end{enumerate}
\end{exercise}

\subsection{\texorpdfstring{$R^2$}{R-squared}}
Assume from now on that $X$ contains a column of ones (an intercept). Then $\bm 1\T\hat{\bm\varepsilon}=0$ by
\eqref{ch05:eq-normal}: residuals sum to zero and $\bar{\hat y}=\bar y$. Subtracting $\bar y\bm 1$ (which lies in
$\operatorname{col}(X)$) from both sides of Pythagoras gives the \emph{analysis of variance}
\[
  \underbrace{\sum_i(y_i-\bar y)^2}_{\text{TSS}}=\underbrace{\sum_i(\hat y_i-\bar y)^2}_{\text{ESS}}
  +\underbrace{\sum_i\hat\varepsilon_i^2}_{\text{RSS}} .
\]

\begin{definition}[Coefficient of determination]\label{ch05:def-r2}
$R^2=\text{ESS}/\text{TSS}=1-\text{RSS}/\text{TSS}$ (``multiple $R$-squared''), and the \emph{adjusted}
$R^2_{\text{adj}}=1-\dfrac{\text{RSS}/(n-p)}{\text{TSS}/(n-1)}$, which penalises the number of parameters.
\end{definition}

\begin{proposition}\label{ch05:prop-r2corr}
With an intercept, $R^2=\Corr(\bm y,\hat{\bm y})^2$, the squared sample correlation between observed and fitted
values \ts{24}{8}{59:31}.
\end{proposition}
\begin{proof}
Let $\tilde{\bm y}=\bm y-\bar y\bm 1$ and $\tilde{\hat{\bm y}}=\hat{\bm y}-\bar y\bm 1$. Since $\hat{\bm\varepsilon}\perp\hat{\bm y}$ and
$\hat{\bm\varepsilon}\perp\bm 1$, $\tilde{\bm y}\T\tilde{\hat{\bm y}}=(\tilde{\hat{\bm y}}+\hat{\bm\varepsilon})\T\tilde{\hat{\bm y}}=\|\tilde{\hat{\bm y}}\|^2$, so
$\Corr^2=\|\tilde{\hat{\bm y}}\|^4/(\|\tilde{\bm y}\|^2\|\tilde{\hat{\bm y}}\|^2)=\text{ESS}/\text{TSS}$.
\end{proof}
For simple regression ($p=2$, one slope) this is the familiar squared correlation of $x$ and $y$, and
$\hat\beta_1=\sum_i(x_i-\bar x)(y_i-\bar y)/\sum_i(x_i-\bar x)^2=\widehat{\Cov}(x,y)/\widehat{\Var}(x)$, the beta
of \cref{def:alphabeta}. It is worth calibrating the eye: the daily CAPM regression of GE on the market in
\cref{ch05:capm} has $R^2=0.30$, a clearly visible but very noisy linear cloud \ts{24}{11}{59:45}.

\begin{exercise}[not from the course: $R^2$ with orthogonal regressors]\label{ch05:ex-r2orth}
Let $\bm z_1,\dots,\bm z_m$ be centred and mutually orthogonal, and regress $\bm y$ on $\bm 1,\bm z_1,\dots,\bm z_m$.
\begin{enumerate}[(a)]
  \item Show that the multiple
      $R^2$ equals the sum of the $R^2$'s of the $m$ simple regressions of $\bm y$ on $\bm 1,\bm z_j$, and that each slope is the same in the
      simple and in the multiple regression.
  \item Check this on the ETF numbers ($0.545+0.006+0.079+0.012$ vs $0.642$) and
      explain why the order of the principal components by variance says nothing about their $R^2$.
\end{enumerate}
\end{exercise}

\subsection{The QR decomposition}\label{ch05:sec-qr}
Numerically one never inverts $X\T X$. The 2013 lecture \ts{13}{6}{1:08:26} uses the QR decomposition
\[
  X=QR,\qquad Q\in\R^{n\times p},\ Q\T Q=I_p,\qquad R\in\R^{p\times p}\ \text{upper triangular},
\]
obtained by Gram--Schmidt on the columns of $X$ \ts{13}{6}{1:09:26}. Column by column,
$X_{[1]}=Q_{[1]}r_{11}$, $X_{[2]}=Q_{[1]}r_{12}+Q_{[2]}r_{22}$, \dots, so
\[\begin{gathered}
  r_{11}=\|X_{[1]}\|,\quad Q_{[1]}=\frac{X_{[1]}}{r_{11}};\\
  r_{12}=Q_{[1]}\T X_{[2]},\quad r_{22}=\|X_{[2]}-r_{12}Q_{[1]}\|,\quad Q_{[2]}=\frac{X_{[2]}-r_{12}Q_{[1]}}{r_{22}};\ \dots
  \end{gathered}\]
(each column is orthogonalised against the previous ones) \ts{13}{6}{1:10:32}. Plugging $X=QR$ into the OLS
formulas \ts{13}{6}{1:11:35}:
\[
  \hat{\bm\beta}=R^{-1}Q\T\bm y,\qquad \Cov(\hat{\bm\beta})=\sigma^2(X\T X)^{-1}=\sigma^2R^{-1}(R^{-1})\T,\qquad H=QQ\T .
\]
$\hat{\bm\beta}$ is computed by back-substitution in $R\hat{\bm\beta}=Q\T\bm y$. The advantage (not stressed in
class): the condition number of $X\T X$ is the square of that of $X$, so forming $X\T X$ loses twice as many
digits when the regressors are nearly collinear. R's \texttt{lm()} works through a QR decomposition (the
\texttt{qr} component of a fitted \texttt{lm} object). The SVD route of 2013 PS1 B-4 (\cref{ch:linalg}) is
the most stable of all and is the one used for ridge and PC regression in \cref{ch05:sec-shrink}.

\subsection{Reparametrisation and standardised covariates}\label{ch05:sec-standard}
If $X'=XG$ with $G\in\R^{p\times p}$ invertible, then $\operatorname{col}(X')=\operatorname{col}(X)$, so $H$, $\hat{\bm y}$ and
$\hat{\bm\varepsilon}$ are unchanged and the new parameter is $\bm\beta'=G^{-1}\bm\beta$ (2013 PS3, \cref{ch05:ex-ps3-1de}).
The important special case is \emph{standardisation} \ts{24}{8}{54:05}: replace each covariate by its
$z$-score $(X_{[j]}-\bar x_j\bm 1)/s_j$. Then
\begin{itemize}
  \item fitted values, residuals, $R^2$, $F$ statistic, and the $t$ statistics and $p$-values of the slopes are
        identical \ts{24}{8}{57:24};
  \item the slope coefficient becomes $s_j\beta_j$: the effect on $y$ of a \emph{one-standard-deviation} move
        of $x_j$, which makes coefficients of variables in different units comparable \ts{24}{8}{58:27};
  \item the intercept changes meaning: it becomes $\bar y$, the prediction at the average covariate.
\end{itemize}
In the prostate example of \cref{ch05:sec-diag} the coefficient of \texttt{lcavol} goes from $0.5435$ to
$0.6080$ (its training-sample standard deviation is $1.119$) and that of \texttt{svi} from $0.877$ to
$0.349$, while the $t$ values stay $5.206$ and $2.949$.

% ------------------------------------------------------------------
\section{The Gauss--Markov theorem}\label{ch05:sec-gm}

\begin{definition}[Gauss--Markov assumptions]
The data $(\bm y,X)$ follow a linear model satisfying the Gauss--Markov assumptions if $\bm y$ is an
observation of a random vector $\bm Y$ with \ts{24}{8}{1:11:00} \ts{13}{6}{45:22}
\[
  \E[\bm Y\mid X,\bm\beta]=X\bm\beta,\qquad \Cov(\bm Y\mid X,\bm\beta)=\sigma^2I_n\quad(\sigma^2>0),
\]
equivalently $\bm Y=X\bm\beta+\bm\varepsilon$ with $\E[\bm\varepsilon]=\bm 0$, $\Cov(\bm\varepsilon)=\sigma^2I_n$. Everything is conditional
on $X$, which is treated as fixed.
\end{definition}

For a random vector $\bm Y$ with mean $\bm\mu$ and covariance $\Sigma$ and a constant matrix $A$,
\[
  \E[A\bm Y]=A\bm\mu,\qquad \Cov(A\bm Y)=A\,\Sigma\,A\T
\]
\ts{13}{6}{1:12:37}. With $A=(X\T X)^{-1}X\T$ this gives, under Gauss--Markov,
\begin{equation}\label{ch05:eq-covbeta}
  \E[\hat{\bm\beta}]=(X\T X)^{-1}X\T X\bm\beta=\bm\beta,\qquad
  \Cov(\hat{\bm\beta})=\sigma^2(X\T X)^{-1}X\T X(X\T X)^{-1}=\sigma^2(X\T X)^{-1}.
\end{equation}

\begin{theorem}[Gauss--Markov]\label{ch05:thm-gm}
For known constants $c_1,\dots,c_{p+1}$ let $\theta=c_1\beta_1+\dots+c_p\beta_p+c_{p+1}$ and
$\hat\theta=c_1\hat\beta_1+\dots+c_p\hat\beta_p+c_{p+1}$ with the OLS estimates. Then $\hat\theta$ is
\begin{enumerate}[(i)]
  \item unbiased;
  \item linear in $\bm y$, $\hat\theta=\sum_ib_iy_i$ with constants $b_i$ depending on $X$ only.
\end{enumerate}
Under the Gauss--Markov assumptions it has the smallest variance among \emph{all} linear unbiased estimators
of $\theta$: it is the \emph{best linear unbiased estimator} (BLUE) \ts{24}{8}{1:12:05} \ts{13}{6}{46:25}.
\end{theorem}
\begin{proof}
(2013 notes) \ts{13}{6}{51:42}. Without loss of generality $c_{p+1}=0$; write $\bm c=(c_1,\dots,c_p)\T$. Then
$\hat\theta=\bm c\T\hat{\bm\beta}=\bm c\T(X\T X)^{-1}X\T\bm y\equiv\bm d\T\bm y$ with $\bm d=X(X\T X)^{-1}\bm c\in\operatorname{col}(X)$, and
$\E[\hat\theta]=\bm c\T\bm\beta$ by \eqref{ch05:eq-covbeta}. Let $\tilde\theta=\bm b\T\bm y$ be any linear unbiased estimator and
$\bm f=\bm b-\bm d$, so $\tilde\theta=\hat\theta+\bm f\T\bm y$. Unbiasedness of both gives
$0=\E[\bm f\T\bm y]=\bm f\T X\bm\beta$ for \emph{all} $\bm\beta\in\R^p$, hence $X\T\bm f=\bm 0$: $\bm f\perp\operatorname{col}(X)$, in particular
$\bm f\perp\bm d$. Therefore
\[
  \Var(\tilde\theta)=\Var(\bm d\T\bm y)+\Var(\bm f\T\bm y)+2\Cov(\bm d\T\bm y,\bm f\T\bm y)
  =\Var(\hat\theta)+\sigma^2\|\bm f\|^2+2\sigma^2\bm d\T\bm f\ \ge\ \Var(\hat\theta),
\]
with equality iff $\bm f=\bm 0$ \ts{13}{6}{53:43}.
\end{proof}

\begin{coursenote}[Erratum in the 2013 proof]
Slide 17 of \file{lecnote6} writes $0=\E[\bm f\T\bm y]=\bm d\T\E[\bm y]=\bm f\T(X\bm\beta)$; the middle term should be
$\bm f\T\E[\bm y]$. The key point, easy to miss, is that unbiasedness must hold \emph{for every} $\bm\beta$, which is
what forces $\bm f\perp\operatorname{col}(X)$.
\end{coursenote}

\paragraph{Remarks.}
\begin{itemize}
  \item The theorem is very general \ts{24}{8}{1:13:07}: with $\bm c=\bm x_0$ it says that $\bm x_0\T\hat{\bm\beta}$ is the best
        linear unbiased estimate of the mean response at a new point $\bm x_0$; with $\bm c=\bm x_1-\bm x_2$, of the
        difference of mean responses of two cases. Since it holds for every $\bm c$, it is equivalent to the
        matrix statement $\Cov(\tilde{\bm\beta})-\Cov(\hat{\bm\beta})\succeq0$ for every linear unbiased $\tilde{\bm\beta}$.
  \item No normality is needed. Econometrics texts call this the BLUE property of OLS \ts{13}{6}{54:45}.
  \item ``Best'' only among \emph{linear unbiased} estimators. A biased estimator (ridge,
        \cref{ch05:sec-shrink}) or a non-linear one (robust M-estimators, \cref{ch05:sec-mle}) can have a smaller
        mean squared error, especially with heavy-tailed errors.
\end{itemize}

\paragraph{Design: where to put the $x$'s} \ts{24}{8}{27:37}. $\Var(\hat\beta_j)=\sigma^2C_{jj}$ with
$C_{jj}=[(X\T X)^{-1}]_{jj}$. If we control $X$ (an experiment), we should make $C_{jj}$ small. If $X\T X$ is
diagonal, $C_{jj}=1/\|X_{[j]}\|^2$: the longer the column, the more precise the estimate \ts{24}{8}{29:50}. For a
straight line $y=\beta_0+\beta_1x$,
\[
  \Var(\hat\beta_1)=\frac{\sigma^2}{\sum_i(x_i-\bar x)^2}:
\]
points far from $\bar x$ pin the slope down \ts{24}{8}{30:55}. In general the precision is governed by the
eigenvalues of $X\T X$ (optimal experimental design) \ts{24}{8}{31:59}.

\begin{intuition}[Lever arms]
A physicist measuring a spring constant spreads the loads over the widest possible range: an error
$\delta y$ at a point at distance $d$ from the centroid tilts the line by $\delta y/d$. The same ``lever arm'' will
reappear as the \emph{leverage} $H_{ii}$ of a case (\cref{ch05:sec-diag}): the points that pin the line down are
also the ones that can pull it away.
\end{intuition}

\begin{exercise}[not from the course: optimal design of a line fit]
For $y=\beta_0+\beta_1x+\varepsilon$ under Gauss--Markov with $n$ even and $x_i\in[-1,1]$ at your choice:
\begin{enumerate}[(a)]
  \item derive
      $\hat\beta_1=S_{xy}/S_{xx}$ and $\Var(\hat\beta_1)=\sigma^2/S_{xx}$ with $S_{xx}=\sum_i(x_i-\bar x)^2$;
  \item find the design minimising
      $\Var(\hat\beta_1)$ and compute the leverages $H_{ii}$ for it;
  \item explain why this design cannot detect curvature, and how the
      curvature test of \texttt{residualPlots()} would behave.
\end{enumerate}
\end{exercise}

% ------------------------------------------------------------------
\section{Generalised least squares}\label{ch05:sec-gls}
Suppose now $\bm Y=X\bm\beta+\bm\varepsilon$ with $\E[\bm\varepsilon]=\bm 0$ and $\Cov(\bm\varepsilon)=\sigma^2\Sigma$, where $\sigma^2$ is an unknown scale
and $\Sigma$ is a \emph{known} $n\times n$ positive definite matrix giving the relative variances and the
correlations of the observations \ts{24}{8}{1:13:07} \ts{13}{6}{54:45}. By the spectral theorem
(\cref{thm:spectral}) $\Sigma=U\Lambda U\T$ has the symmetric positive definite inverse square root
$\Sigma^{-1/2}=U\Lambda^{-1/2}U\T$. Transform the data \ts{24}{8}{1:15:23} \ts{13}{6}{55:47}:
\[\begin{gathered}
  \bm Y^*=\Sigma^{-1/2}\bm Y,\quad X^*=\Sigma^{-1/2}X,\quad \bm\varepsilon^*=\Sigma^{-1/2}\bm\varepsilon:\\
  \bm Y^*=X^*\bm\beta+\bm\varepsilon^*,\quad \Cov(\bm\varepsilon^*)=\sigma^2\Sigma^{-1/2}\Sigma\Sigma^{-1/2}=\sigma^2I_n .
\end{gathered}\]
The transformed model satisfies Gauss--Markov with the \emph{same} $\bm\beta$, so by \cref{ch05:thm-gm} the BLUE is
OLS on the starred data:
\begin{equation}\label{ch05:eq-gls}
  \hat{\bm\beta}_{\text{GLS}}=\bigl(X^{*\mathsf T}X^*\bigr)^{-1}X^{*\mathsf T}\bm Y^*=\bigl(X\T\Sigma^{-1}X\bigr)^{-1}X\T\Sigma^{-1}\bm Y,
  \qquad \Cov(\hat{\bm\beta}_{\text{GLS}})=\sigma^2\bigl(X\T\Sigma^{-1}X\bigr)^{-1}.
\end{equation}
Equivalently $\hat{\bm\beta}_{\text{GLS}}$ minimises the Mahalanobis-type criterion $(\bm y-X\bm\beta)\T\Sigma^{-1}(\bm y-X\bm\beta)$. If
$\Sigma=\diag(\sigma_1^2,\dots,\sigma_n^2)$ this is \emph{weighted least squares} with weights $1/\sigma_i^2$: noisy cases are
down-weighted \ts{24}{8}{1:16:24} \ts{13}{6}{57:55}. The main source of non-diagonal $\Sigma$ in finance is time
dependence: errors at nearby dates are more correlated than errors far apart \ts{24}{8}{1:14:15}; e.g.\ AR(1)
errors have $\Sigma_{ij}\propto\rho^{|i-j|}$ (\cref{ch:timeseries}).

\begin{intuition}[Whitening]
$\chi^2=\bm r\T C^{-1}\bm r$ with the covariance matrix $C$ of correlated measurements is exactly the GLS criterion.
Multiplying by $\Sigma^{-1/2}$ ``whitens'' the noise: in the new coordinates the noise is isotropic and the plain
Euclidean projection of \cref{ch05:fig-geometry} is optimal again.
\end{intuition}

\begin{remark}[Not discussed in class]
If the errors are correlated but OLS is used anyway, $\hat{\bm\beta}$ is still unbiased, but it is no longer best and its
true covariance is the ``sandwich'' $\sigma^2(X\T X)^{-1}X\T\Sigma X(X\T X)^{-1}$, not $\sigma^2(X\T X)^{-1}$: the standard
errors printed by \texttt{lm()} are wrong. In practice $\Sigma$ is unknown; it is estimated (feasible GLS) or the
standard errors are corrected (heteroscedasticity- and autocorrelation-consistent errors).
\end{remark}

% ------------------------------------------------------------------
\section{The normal linear model: distribution theory}\label{ch05:sec-normal}

\subsection{Multivariate normal vectors via MGFs}
For an $n$-variate random vector $\bm Y$ the moment-generating function is $M_{\bm Y}(\bm t)=\E[e^{\bm t\T\bm Y}]$,
$\bm t\in\R^n$ \ts{24}{8}{21:10}. In the normal linear model the $Y_i$ are independent $N(\mu_i,\sigma^2)$ with
$\bm\mu=X\bm\beta$ (independent but not identically distributed \ts{13}{6}{57:55}), so the expectation of the product
factorises \ts{24}{8}{22:13} \ts{13}{6}{1:02:09}:
\[
  M_{\bm Y}(\bm t)=\prod_{i=1}^n\E[e^{t_iY_i}]=\prod_{i=1}^ne^{t_i\mu_i+\frac12t_i^2\sigma^2}
  =\exp\Bigl(\bm t\T\bm\mu+\tfrac12\bm t\T\Sigma\bm t\Bigr),\qquad \Sigma=\sigma^2I_n .
\]

\begin{definition}[Multivariate normal]\label{ch05:def-mvn}
$\bm Y\sim N_n(\bm\mu,\Sigma)$, with $\bm\mu\in\R^n$ and $\Sigma$ symmetric positive semi-definite, if
$M_{\bm Y}(\bm t)=\exp(\bm t\T\bm\mu+\frac12\bm t\T\Sigma\bm t)$ for all $\bm t$. Then $\E[\bm Y]=\bm\mu$ and $\Cov\bm Y=\Sigma$; if
$\Sigma\succ0$ the density is $(2\pi)^{-n/2}(\det\Sigma)^{-1/2}\exp\bigl(-\frac12(\bm y-\bm\mu)\T\Sigma^{-1}(\bm y-\bm\mu)\bigr)$. Since an MGF
finite around $\bm 0$ determines the distribution, recognising this form is enough to identify a normal vector.
\end{definition}

\begin{coursenote}[Vocabulary]
In the 2024 L8 video Kempthorne often says ``multinomial'' \ts{24}{8}{20:07}: he means \emph{multinormal},
i.e.\ multivariate normal. Nothing in this chapter involves the multinomial distribution.
\end{coursenote}

\begin{proposition}[Linear maps of normal vectors]\label{ch05:prop-linmap}
If $\bm Y\sim N_n(\bm\mu,\Sigma)$ and $A\in\R^{m\times n}$ then $\bm Z=A\bm Y\sim N_m(A\bm\mu,A\Sigma A\T)$ \ts{24}{8}{31:59}
\ts{13}{6}{1:14:49}. In particular every sub-vector, and every component, is normal; and two jointly normal
sub-vectors with zero cross-covariance are independent.
\end{proposition}
\begin{proof}
$M_{\bm Z}(\bm\tau)=\E[e^{\bm\tau\T A\bm Y}]=M_{\bm Y}(A\T\bm\tau)=\exp\bigl(\bm\tau\T A\bm\mu+\frac12\bm\tau\T A\Sigma A\T\bm\tau\bigr)$. For a component take
$A=\bm e_j\T$ \ts{13}{6}{1:07:25}. If $\bm Z=(\bm Z_1,\bm Z_2)$ is normal with $\Cov(\bm Z_1,\bm Z_2)=0$, the quadratic form in the
exponent splits and the joint MGF is the product $M_{\bm Z_1}(\bm\tau_1)M_{\bm Z_2}(\bm\tau_2)$, which characterises
independence \ts{24}{8}{37:20}.
\end{proof}
Applying this with $A=(X\T X)^{-1}X\T$ \ts{24}{8}{24:21}: $\bm t=A\T\bm\tau=X(X\T X)^{-1}\bm\tau$ gives $\bm t\T\bm\mu=\bm\tau\T\bm\beta$ and
$\bm t\T(\sigma^2I)\bm t=\bm\tau\T\sigma^2(X\T X)^{-1}\bm\tau$, so $\hat{\bm\beta}\sim N_p(\bm\beta,\sigma^2(X\T X)^{-1})$ and
$\hat\beta_j\sim N(\beta_j,\sigma^2C_{jj})$ \ts{24}{8}{26:31} \ts{13}{6}{1:04:19}.

\begin{coursenote}[Errata in the MGF slides]
\begin{itemize}
  \item Slide 17 of \file{lec06_1}: the MGF of $\hat\beta_1$ is obtained with $\bm\tau=t(1,0,\dots,0)\T$, not $t(1,0,\dots,1)\T$.
  \item Slides 14--16 write the exponent as $\bm t\T\bm u+\dots$: $\bm u$ is the mean vector $\bm\mu$.
  \item Slide 18 (``Theorem*'') assumes $\rank A=m\le n$. The result holds for every $A$ if degenerate normals
        (singular $\Sigma$, no density) are allowed, as in \cref{ch05:def-mvn}; this is needed right away, because the
        stacked matrix $\bigl(\begin{smallmatrix}(X\T X)^{-1}X\T\\ I-H\end{smallmatrix}\bigr)$ used for independence has $n+p$
        rows but rank $n$, and $\Cov(\hat{\bm\varepsilon})$ is singular.
\end{itemize}
\end{coursenote}

\subsection{The fundamental theorem}
\begin{definition}[$\chi^2$, $t$ and $F$ distributions]
If $Z_1,\dots,Z_\nu$ are i.i.d.\ $N(0,1)$, $\sum_kZ_k^2\sim\chi^2_\nu$ \ts{13}{6}{1:16:57}. If $Z\sim N(0,1)$ and
$V\sim\chi^2_\nu$ are independent, $Z/\sqrt{V/\nu}\sim t_\nu$ (Student) \ts{24}{8}{38:21}. If $V_1\sim\chi^2_{\nu_1}$ and
$V_2\sim\chi^2_{\nu_2}$ are independent, $(V_1/\nu_1)/(V_2/\nu_2)\sim F_{\nu_1,\nu_2}$ \ts{24}{8}{45:38}. Note $t_\nu^2\sim F_{1,\nu}$.
\end{definition}

\begin{theorem}[Normal linear regression]\label{ch05:thm-normal}
Let $\bm y=X\bm\beta+\bm\varepsilon$ with $X\in\R^{n\times p}$ of rank $p$ and $\bm\varepsilon\sim N_n(\bm 0,\sigma^2I_n)$. Then
\ts{24}{8}{33:00} \ts{13}{6}{1:15:54}:
\begin{enumerate}[(i)]
  \item $\hat{\bm\beta}=(X\T X)^{-1}X\T\bm y\sim N_p\bigl(\bm\beta,\sigma^2(X\T X)^{-1}\bigr)$;
  \item $\hat{\bm\varepsilon}=(I_n-H)\bm y\sim N_n\bigl(\bm 0,\sigma^2(I_n-H)\bigr)$, a degenerate normal: $X\T\hat{\bm\varepsilon}=\bm 0$ identically
        \ts{24}{8}{34:06};
  \item $\hat{\bm\beta}$ and $\hat{\bm\varepsilon}$ are independent;
  \item $\text{RSS}=\hat{\bm\varepsilon}\T\hat{\bm\varepsilon}\sim\sigma^2\chi^2_{n-p}$; hence $\hat\sigma^2=\text{RSS}/(n-p)$ is unbiased;
  \item for each $j$,
        \begin{equation}\label{ch05:eq-tstat}
          \hat t_j=\frac{\hat\beta_j-\beta_j}{\hat\sigma\sqrt{C_{jj}}}\sim t_{n-p},\qquad C_{jj}=\bigl[(X\T X)^{-1}\bigr]_{jj}.
        \end{equation}
\end{enumerate}
\end{theorem}
\begin{proof}
(i) was shown above.

(ii) $\Cov\bigl((I-H)\bm y\bigr)=\sigma^2(I-H)(I-H)\T=\sigma^2(I-H)$ and $(I-H)X\bm\beta=\bm 0$.

(iii) $\hat{\bm\beta}$ and $\hat{\bm\varepsilon}$ are jointly normal (both linear in $\bm y$) with
$\Cov(\hat{\bm\beta},\hat{\bm\varepsilon})=(X\T X)^{-1}X\T(\sigma^2I)(I-H)=\sigma^2(X\T X)^{-1}(X\T-X\T H)=0$ because $X\T H=X\T$; apply
\cref{ch05:prop-linmap} (this is the ``joint MGF = product of MGFs'' argument of the 2024 slides).

(iv) Following the 2013 notes \ts{13}{6}{1:16:57}: let $X=QR$ and complete the columns of $Q$ to an orthonormal
basis of $\R^n$ by $W\in\R^{n\times(n-p)}$ (continue Gram--Schmidt on $[X\,|\,I_n]$). Then
$A=\bigl(\begin{smallmatrix}Q\T\\ W\T\end{smallmatrix}\bigr)$ is orthogonal and
\[
  \bm z=A\bm y=\begin{pmatrix}\bm z_Q\\ \bm z_W\end{pmatrix}\sim N_n\!\left(\begin{pmatrix}Q\T QR\bm\beta\\ W\T QR\bm\beta\end{pmatrix},\sigma^2AA\T\right)
  =N_n\!\left(\begin{pmatrix}R\bm\beta\\ \bm 0_{n-p}\end{pmatrix},\sigma^2I_n\right).
\]
So $\bm z_Q\sim N_p(R\bm\beta,\sigma^2I_p)$ and $\bm z_W\sim N_{n-p}(\bm 0,\sigma^2I_{n-p})$ are independent. Now
$\hat{\bm\beta}=R^{-1}Q\T\bm y=R^{-1}\bm z_Q$, and since $QQ\T+WW\T=A\T A=I_n$,
$\hat{\bm\varepsilon}=(I-QQ\T)\bm y=WW\T\bm y=W\bm z_W$. Hence $\text{RSS}=\|W\bm z_W\|^2=\|\bm z_W\|^2$ is $\sigma^2$ times a sum of $n-p$
independent squared standard normals. (This also re-proves (iii): $\hat{\bm\beta}$ and $\hat{\bm\varepsilon}$ are functions of the
independent blocks $\bm z_Q$ and $\bm z_W$.)

(v) $(\hat\beta_j-\beta_j)/(\sigma\sqrt{C_{jj}})\sim N(0,1)$ by (i), independent of $\text{RSS}/\sigma^2\sim\chi^2_{n-p}$ by (iii)--(iv);
dividing, $\sigma$ cancels and the ratio is \eqref{ch05:eq-tstat}.
\end{proof}

The unbiasedness of $\hat\sigma^2$ does not need normality \ts{24}{8}{36:15}: under Gauss--Markov,
$\E[\hat{\bm\varepsilon}\T\hat{\bm\varepsilon}]=\tr\Cov(\hat{\bm\varepsilon})=\sigma^2\tr(I_n-H)=\sigma^2(n-p)$ by \cref{ch05:prop-hat}(iii). The
$n-p$ appears because the residuals satisfy $p$ exact linear constraints $X\T\hat{\bm\varepsilon}=\bm 0$: there are not $n$
independent error terms \ts{24}{8}{35:14}.

\begin{coursenote}[Erratum: the missing square root]
Slide 20 of \file{lec06_1} and slide 34 of \file{lecnote6} write the denominator of $\hat t_j$ as
$\hat\sigma\,C_{j,j}$. It must be $\hat\sigma\sqrt{C_{jj}}$, the estimated standard deviation of $\hat\beta_j$; Kempthorne
corrects it at the start of L11 \ts{24}{11}{3:16}. Also, slide 38 of \file{lecnote6} writes
$\hat{\bm\beta}=(X\T X)^{-1}X\bm y$ for $(X\T X)^{-1}X\T\bm y$.
\end{coursenote}

\begin{intuition}[Why Student's $t$ has fat tails]
In $\hat t_j$ the unknown $\sigma$ is replaced by the random $\hat\sigma$. When $\hat\sigma$ happens to be small, $|\hat t_j|$ is
large; the denominator fluctuates around $1$ (in units of $\sigma$) and this mixing over scales fattens the tails
\ts{24}{8}{40:28}. W.~S.~Gosset, a statistician at the Guinness brewery, noticed that standardised means of
quality samples of size four were far more variable than $N(0,1)$ predicted; the brewery did not let him
publish under his name, hence ``Student'' \ts{24}{8}{41:29}. For $n-p\gtrsim30$ the
difference from the normal is small; with $n-p=63$ (prostate example) the two-sided 5\% critical value is
$2.00$ instead of $1.96$.
\end{intuition}

\begin{intuition}[Degrees of freedom, geometrically]
$\text{RSS}/\sigma^2$ is the squared length of a standard Gaussian vector projected on an $(n-p)$-dimensional subspace:
fitting $p$ parameters ``uses up'' $p$ dimensions. The familiar $1/(n-1)$ in the sample variance is the case
$p=1$ (projection on $\bm 1$), and $\chi^2/\mathrm{ndf}$ in a fit is $\hat\sigma^2/\sigma^2$.
\end{intuition}

A $(1-\alpha)$ confidence interval for $\beta_j$ follows immediately:
$\hat\beta_j\pm t_{n-p,1-\alpha/2}\,\hat\sigma\sqrt{C_{jj}}$, where $\hat\sigma\sqrt{C_{jj}}$ is the \emph{standard error}
$\mathrm{se}(\hat\beta_j)$ printed by R.

% ------------------------------------------------------------------
\section{Maximum likelihood and M-estimation}\label{ch05:sec-mle}

\subsection{Maximum likelihood in the normal model}
The \emph{likelihood} is the joint density of the data seen as a function of the unknown parameters,
$L(\bm\beta,\sigma^2)=p(\bm y\mid X,\bm\beta,\sigma^2)$; maximum likelihood (ML) estimates make the observed data ``most
likely'' \ts{24}{11}{10:53} \ts{13}{6}{1:20:05}. For the normal linear model
\[
  L(\bm\beta,\sigma^2)=\prod_{i=1}^n\frac{1}{\sqrt{2\pi\sigma^2}}e^{-\frac{(y_i-\bm x_i\T\bm\beta)^2}{2\sigma^2}},\qquad
  \log L=-\frac n2\log(2\pi\sigma^2)-\frac{1}{2\sigma^2}Q(\bm\beta)
\]
\ts{24}{11}{13:02}. Maximise in two steps: $\bm\beta$ enters only through $-Q(\bm\beta)$, so $\hat{\bm\beta}_{\text{ML}}=\hat{\bm\beta}_{\text{OLS}}$
for every $\sigma^2$; then $\partial\log L/\partial\sigma^2=-\frac{n}{2\sigma^2}+\frac{Q(\hat{\bm\beta})}{2\sigma^4}=0$ gives
\[
  \hat\sigma^2_{\text{ML}}=\frac{\text{RSS}}{n},\qquad \E[\hat\sigma^2]_{\text{ML}}=\frac{n-p}{n}\sigma^2\quad\text{(biased)}
\]
\ts{24}{11}{14:08}. So least squares is also ``doing something decent'' from the likelihood point of view when the
errors are Gaussian \ts{13}{6}{1:21:11}. ML estimators are asymptotically optimal (smallest variance as
$n\to\infty$) for any correctly specified model \ts{24}{11}{11:58}: if the error density were known and non-Gaussian,
ML would use it.

\subsection{Generalised M-estimators}
\begin{definition}[Generalised M-estimator]\label{ch05:def-mest}
For data $(\bm y,X)$ and the model $y_i=\bm x_i\T\bm\beta+\varepsilon_i$, choose $\hat{\bm\beta}$ to minimise
$Q(\bm\beta)=\sum_{i=1}^nh(y_i,\bm x_i,\bm\beta,\sigma^2)$ \ts{24}{11}{15:14}. The function $h$ defines the estimator:
\begin{enumerate}[(i)]
  \item least squares: $h=(y_i-\bm x_i\T\bm\beta)^2$;
  \item least absolute deviations (``MAD''): $h=|y_i-\bm x_i\T\bm\beta|$;
  \item maximum likelihood: $h=-\log p(y_i\mid\bm\beta,\bm x_i,\sigma^2)$ (so Laplace errors give LAD);
  \item robust M-estimator: $h=\chi(y_i-\bm x_i\T\bm\beta)$ with $\chi$ even and increasing on $(0,\infty)$;
  \item quantile estimator, for fixed $\tau\in(0,1)$: $h=\rho_\tau(y_i-\bm x_i\T\bm\beta)$ with the \emph{check loss}
        \[
          \rho_\tau(u)=\begin{cases}\tau\,|u|,&u\ge0,\\(1-\tau)\,|u|,&u<0.\end{cases}
        \]
\end{enumerate}
\end{definition}

\begin{figure}[ht]
\centering
\begin{tikzpicture}
\begin{axis}[width=0.62\linewidth,height=5.2cm,xmin=-2.5,xmax=2.5,ymin=0,ymax=2.6,
  axis lines=middle,xlabel={$u$},ylabel={$h(u)$},samples=101,
  legend style={font=\scriptsize,at={(1.02,1)},anchor=north west,draw=none},
  every axis x label/.style={at={(current axis.right of origin)},anchor=west},
  every axis y label/.style={at={(current axis.above origin)},anchor=south}]
  \addplot[thick,black!60,dashed,domain=-2.3:2.3]{x^2/2};
  \addlegendentry{$u^2/2$ (least squares)}
  \addplot[thick,defcol,domain=-2.4:2.4]{abs(x)};
  \addlegendentry{$|u|$ (LAD)}
  \addplot[thick,intcol,domain=-2.4:2.4]{(abs(x)<=1)*x^2/2+(abs(x)>1)*(abs(x)-0.5)};
  \addlegendentry{Huber, $c=1$}
  \addplot[very thick,thmcol,domain=-2.4:2.4]{(x>=0)*0.9*x+(x<0)*(-0.1)*x};
  \addlegendentry{$\rho_{0.9}$ (quantile)}
\end{axis}
\end{tikzpicture}
\caption{Loss functions of M-estimators. Quadratic loss lets large residuals dominate; LAD and Huber grow
linearly; the check loss $\rho_{0.9}$ penalises positive residuals nine times more than negative ones, so the fit
is pushed up to the 90th percentile.}\label{ch05:fig-losses}
\end{figure}

Robust estimators were designed for the contaminated-normal situation: normal errors most of the time, with a
small fraction of gross errors \ts{24}{11}{17:30}. When the error distribution is unknown they protect the fit
from a few extreme cases, which in finance means crash days. A standard example (not detailed in class) is
Huber's loss, quadratic for $|u|\le c$ and linear beyond.

\begin{proposition}[The check loss estimates quantiles]\label{ch05:prop-quantile}
Let $y_1,\dots,y_n\in\R$ and $f(c)=\sum_i\rho_\tau(y_i-c)$. Then $f$ is convex and piecewise linear, and $c$ minimises $f$ iff
\[
  \#\{i:y_i<c\}\le\tau n\le\#\{i:y_i\le c\},
\]
i.e.\ iff $c$ is a sample $\tau$-quantile. For $\tau=\frac12$ this is the median, and $\rho_{1/2}=\frac12|u|$, so
$\tau=\frac12$ gives LAD \ts{24}{11}{18:37}.
\end{proposition}
\begin{proof}
Each term is convex in $c$. The right derivative of $\rho_\tau(y_i-c)$ is $-\tau$ if $y_i>c$ and $1-\tau$ if $y_i\le c$,
so $f'_+(c)=-\tau\,\#\{y_i>c\}+(1-\tau)\,\#\{y_i\le c\}=\#\{y_i\le c\}-\tau n$; similarly $f'_-(c)=\#\{y_i<c\}-\tau n$. A convex
function is minimised exactly where $f'_-(c)\le0\le f'_+(c)$.
\end{proof}
With $c$ replaced by $\bm x_i\T\bm\beta$, the minimiser estimates the conditional $\tau$-quantile of $y$ given $\bm x$
(quantile regression): for $\tau=0.9$, the upper decile of outcomes \ts{24}{11}{19:41}. The objective is convex
(``a bowl with straight edges'') but not differentiable, so there are no normal equations to solve; one moves
downhill using one-sided derivatives, or solves a linear programme \ts{24}{11}{20:41}.

\begin{finance}[Heavy tails in practice]
In the GE CAPM regression of \cref{ch05:capm} the histogram of residuals is compared with two normal fits
\ts{24}{11}{1:00:53}: the ML fit (sample mean and standard deviation) and a robust one, whose scale is
estimated from the interquartile range ($\hat\sigma_{\text{rob}}=\mathrm{IQR}/1.349$, since the IQR of $N(0,\sigma^2)$ is
$1.349\sigma$) \ts{24}{11}{1:03:05}. The ML curve is too wide in the centre and too thin in the tails. A sharper
check uses \emph{fitted percentiles} $u_i=\Phi\bigl((\hat\varepsilon_i-\hat\mu)/\hat\sigma\bigr)$, which should be uniform on $(0,1)$ if the
model is right (probability integral transform) \ts{24}{11}{1:04:09}: with the ML scale they pile up near
$0.5$, with the robust scale they are uniform except in the extreme 1\% on each side
\ts{24}{11}{1:05:11}. Quantile regression is the natural tool for Value-at-Risk (\cref{ch:var}).
\end{finance}

\begin{exercise}[not from the course: quantile regression as maximum likelihood]
Show that minimising $\sum_i\rho_\tau(y_i-\bm x_i\T\bm\beta)$ is maximum likelihood for errors with the \emph{asymmetric Laplace} density
$f(u)=\frac{\tau(1-\tau)}{s}\exp\bigl(-\rho_\tau(u)/s\bigr)$, and check that $f$ is a density whose $\tau$-quantile is $0$. Which estimator of
\cref{ch05:def-mest} is ML for Laplace errors?
\end{exercise}

% ------------------------------------------------------------------
\section{Hypothesis testing}\label{ch05:sec-tests}

\subsection{Individual coefficients}
To test $H_0:\beta_j=b$ against $H_1:\beta_j\ne b$ use $\hat t=(\hat\beta_j-b)/\mathrm{se}(\hat\beta_j)$, which is $t_{n-p}$
under $H_0$ by \eqref{ch05:eq-tstat}; reject when $|\hat t|>c$ for a critical value $c$ \ts{24}{11}{5:26}. The
\emph{$p$-value} is the probability, computed under $H_0$, of a statistic at least as extreme as the one observed,
$P(|T_{n-p}|\ge|\hat t|)$ \ts{24}{8}{49:56}. The coefficient table of R's \texttt{summary(lm(\dots))} has columns
Estimate ($\hat\beta_j$), Std.\ Error ($\hat\sigma\sqrt{C_{jj}}$), $t$ value (test of $b=0$: the ``signal-to-noise
ratio'') and \texttt{Pr(>|t|)} \ts{24}{8}{48:48}. Conventionally $p<0.05$ is ``significant'' and $p<0.01$
``highly significant''. The case $b=0$ asks whether variable $j$ can be dropped from the model
\ts{24}{11}{5:26}.

\begin{example}[Is GE's beta equal to one? (\file{lec08}, section 1)]\label{ch05:ex-gebeta}
Daily excess returns of GE on the market excess return, $n=1676$ days (2017--2023): $\hat\alpha=-0.00054$ (se
$0.00051$, $p=0.285$), $\hat\beta=1.0830$ (se $0.0400$), $R^2=0.304$, $\hat\sigma=0.0207$. Testing $\beta=0$ is pointless
(a stock is obviously correlated with the market, $t=27.1$); the interesting null is $\beta=1$, ``same systematic
risk as the market'' \ts{24}{11}{59:45}:
\[
  \hat t=\frac{1.0830-1}{0.0400}=2.075,\qquad p=P(|T_{1674}|\ge2.075)=0.038,
\]
so GE's beta is (just) significantly above one at the 5\% level \ts{24}{11}{1:07:25}.
\end{example}

\begin{coursenote}[Errata in \file{lec08}, section 1]
On page 3 both hypotheses are labelled $H_0$; the alternative is $H_1:\beta_j\neq0$. On page 4,
``$(\hat\beta_2-\beta_2)\sim N(\beta_2,sd(\hat\beta_2))$'' should read $\hat\beta_2-\beta_2\sim N\bigl(0,sd(\hat\beta_2)^2\bigr)$, i.e.\
$\hat\beta_2\sim N(\beta_2,sd^2)$.
\end{coursenote}

\subsection{Nested models and the \texorpdfstring{$F$}{F} test}
\begin{lemma}[Quadratic forms of normal vectors]\label{ch05:lem-proj}
Let $\bm Y\sim N_n(\bm\mu,\sigma^2I_n)$ and let $P_1,P_2$ be orthogonal projections of ranks $r_1,r_2$ with $P_1P_2=0$ and
$P_1\bm\mu=P_2\bm\mu=\bm 0$. Then $\bm Y\T P_k\bm Y/\sigma^2\sim\chi^2_{r_k}$ ($k=1,2$), and the two quadratic forms are independent.
\end{lemma}
\begin{proof}
Write $P_k=U_kU_k\T$ with $U_k\in\R^{n\times r_k}$ having orthonormal columns (spectral theorem: the eigenvalues of a
projection are $0$ and $1$). Then $U_k\T\bm Y\sim N_{r_k}(U_k\T\bm\mu,\sigma^2I)$ with $U_k\T\bm\mu=U_k\T P_k\bm\mu=\bm 0$, and
$\bm Y\T P_k\bm Y=\|U_k\T\bm Y\|^2$. Moreover $U_1\T U_2=U_1\T(P_1P_2)U_2=0$, so $\Cov(U_1\T\bm Y,U_2\T\bm Y)=\sigma^2U_1\T U_2=0$ and
the two Gaussian vectors are independent (\cref{ch05:prop-linmap}).
\end{proof}

\begin{theorem}[$F$ test of a linear hypothesis]\label{ch05:thm-F}
In the normal linear model let $A\in\R^{q\times p}$ have rank $q$ and consider $H_0:A\bm\beta=\bm 0$. Let RSS be the
residual sum of squares of the unconstrained fit and $\text{RSS}_0$ that of the least-squares fit under the
constraint. Under $H_0$
\[
  F=\frac{(\text{RSS}_0-\text{RSS})/q}{\text{RSS}/(n-p)}\sim F_{q,\,n-p}.
\]
The most common case is $A=[0\ I_q]$: $H_0:\beta_{k+1}=\dots=\beta_p=0$ with $k=p-q$, i.e.\ the \emph{sub-model} using
only the first $k$ columns of $X$ is adequate \ts{24}{8}{43:35} \ts{24}{11}{6:33}.
\end{theorem}
\begin{proof}
The constrained fitted values range over $L_0=\{X\bm\beta:A\bm\beta=\bm 0\}$, a subspace of $\operatorname{col}(X)$ of dimension $p-q$
($X$ is injective). The constrained least-squares fit is the projection $H_0\bm y$ on $L_0$, so
$\text{RSS}_0=\bm y\T(I-H_0)\bm y$ and $\text{RSS}_0-\text{RSS}=\bm y\T(H-H_0)\bm y$. Since $L_0\subset\operatorname{col}(X)$, $HH_0=H_0H=H_0$; hence
$H-H_0$ is the orthogonal projection on $\operatorname{col}(X)\ominus L_0$ (rank $q$) and $(I-H)(H-H_0)=0$. Under $H_0$,
$\bm\mu=X\bm\beta\in L_0$, so $(H-H_0)\bm\mu=(I-H)\bm\mu=\bm 0$. By \cref{ch05:lem-proj}, $(\text{RSS}_0-\text{RSS})/\sigma^2\sim\chi^2_q$ and
$\text{RSS}/\sigma^2\sim\chi^2_{n-p}$ independently, and the ratio of the two over their degrees of freedom is $F_{q,n-p}$.
\end{proof}
Both numerator and denominator estimate $\sigma^2$ when $H_0$ is true \ts{24}{11}{6:33}; under the alternative the
numerator is inflated by $\|(H-H_0)\bm\mu\|^2/q$, so large $F$ is evidence against $H_0$ \ts{24}{8}{44:36}. An affine
hypothesis $A\bm\beta=\bm b$ reduces to this one by subtracting $X\bm\beta_b$ (any $\bm\beta_b$ with $A\bm\beta_b=\bm b$) from $\bm y$.

\begin{coursenote}[Erratum: denominator degrees of freedom]
Slide 20 of \file{lec06_1} gives $df_2=n-k$ for the sub-model $F$ test. It is $df_2=n-p$, the residual degrees of
freedom of the \emph{full} model, as in \file{lec08} (section 2) and in the formula on the same slide.
\end{coursenote}

\begin{proposition}[Wald form; $t^2=F$]\label{ch05:prop-wald}
For $H_0:A\bm\beta=\bm b$,
\[
  \text{RSS}_0-\text{RSS}=(A\hat{\bm\beta}-\bm b)\T\bigl[A(X\T X)^{-1}A\T\bigr]^{-1}(A\hat{\bm\beta}-\bm b).
\]
In particular, for a single coefficient ($A=\bm e_j\T$, $q=1$) $\text{RSS}_0-\text{RSS}=(\hat\beta_j-b)^2/C_{jj}$ and the $F$ statistic
equals $\hat t^{\,2}$ \ts{24}{11}{7:39}.
\end{proposition}
\begin{proof}
Minimising $\|\bm y-X\bm\beta\|^2$ subject to $A\bm\beta=\bm b$ with a Lagrange multiplier gives
$\hat{\bm\beta}_0=\hat{\bm\beta}-(X\T X)^{-1}A\T M^{-1}(A\hat{\bm\beta}-\bm b)$, $M=A(X\T X)^{-1}A\T$. As $X(\hat{\bm\beta}-\hat{\bm\beta}_0)\in\operatorname{col}(X)\perp\hat{\bm\varepsilon}$,
Pythagoras gives $\text{RSS}_0=\text{RSS}+(\hat{\bm\beta}-\hat{\bm\beta}_0)\T X\T X(\hat{\bm\beta}-\hat{\bm\beta}_0)$, and substituting,
$(\hat{\bm\beta}-\hat{\bm\beta}_0)\T X\T X(\hat{\bm\beta}-\hat{\bm\beta}_0)=(A\hat{\bm\beta}-\bm b)\T M^{-1}(A\hat{\bm\beta}-\bm b)$. For $q=1$, dividing by $\hat\sigma^2$ gives
$(\hat\beta_j-b)^2/(\hat\sigma^2C_{jj})=\hat t^{\,2}$.
\end{proof}
R's \texttt{car::linearHypothesis()} implements tests of linear hypotheses in this $F$ form \ts{24}{11}{1:08:29}.
For GE (\cref{ch05:ex-gebeta}) it gives $F=4.304=2.075^2$ ($p=0.038$) for $\beta=1$ and $F=1.144=1.070^2$ ($p=0.285$) for
$\alpha=0$: the same tests as the $t$ statistics.

\paragraph{Overall $F$ and ANOVA.} With an intercept, the test that all $p-1$ slopes vanish has
$\text{RSS}_0=\text{TSS}$, so
\[
  F=\frac{R^2/(p-1)}{(1-R^2)/(n-p)} .
\]
This is the last line of \texttt{summary.lm}: $F=13.4$ on $(8,63)$ degrees of freedom for the prostate data
($R^2=0.630$) and $F=271.9$ on $(4,607)$ for the ETF regression of \cref{ch05:sec-cases} ($R^2=0.642$). Adding
blocks of variables one at a time and recording the drop in RSS gives the sequential \emph{analysis-of-variance
table}, derived through the QR decomposition in 2013 PS3 (\cref{ch05:ex-ps3-3cd}).

\begin{finance}[Statistical versus practical significance; multiple testing]
With thousands of daily observations, tiny effects become ``significant'': in the ETF case study a principal
component explaining $0.6\%$ of the variance of the response has $p=0.052$ \ts{24}{11}{45:48}. Conversely, when the
same test is run on hundreds of stocks (is $\alpha=0$ for each S\&P~500 stock?), about 5\% of them will be
significant at the 5\% level even if every null is true \ts{24}{11}{1:16:47}. A handful of significant alphas is
therefore no evidence against the CAPM, and ``the best of 500 back-tests'' is not a strategy.
\end{finance}

\begin{exercise}[not from the course: $t^2=F$ on real numbers]
\leavevmode
\begin{enumerate}[(a)]
  \item From \cref{ch05:prop-wald}, show directly that for $H_0:\beta_j=b$ the $F$ statistic equals $\hat t^{\,2}$.
  \item Using only the
      GE output of \cref{ch05:ex-gebeta} (estimate, standard error, $n-p=1674$, RSS $=0.71819$), reproduce the RSS $0.72003$ of the
      model constrained by $\beta=1$ printed by \texttt{linearHypothesis}.
\end{enumerate}
\end{exercise}

% ------------------------------------------------------------------
\section{Testing for a change of regime}\label{ch05:sec-chow}
Regressions on time series are fitted over long samples, but the relationship may change: a firm changes business,
a currency changes its peg. \file{lec08} (section 3) shows how to test whether a regression is the same in two
periods \ts{24}{11}{1:09:33}.

\paragraph{Set-up.} Order the $n$ cases in time; period A is $i=1,\dots,n_A$, period B the remaining
$n_B=n-n_A$. Partition $\bm y=(\bm y_A,\bm y_B)$, $X=\bigl(\begin{smallmatrix}X_A\\X_B\end{smallmatrix}\bigr)$ and allow different parameters:
\[
  \bm y_A=X_A\bm\beta+\bm\varepsilon_A,\qquad \bm y_B=X_B\bm\gamma+\bm\varepsilon_B,\qquad \bm\varepsilon_A,\bm\varepsilon_B\ \text{independent},\ N(\bm 0,\sigma^2I),
\]
with $H_0:\bm\gamma=\bm\beta$ (no change) against $H_1:\bm\gamma\neq\bm\beta$ \ts{24}{11}{1:10:37}.

\paragraph{Approach 1: separate regressions.} In block form the unconstrained model is
\[
  \begin{pmatrix}\bm y_A\\ \bm y_B\end{pmatrix}=\begin{pmatrix}X_A&0\\0&X_B\end{pmatrix}\begin{pmatrix}\bm\beta\\ \bm\gamma\end{pmatrix}+\bm\varepsilon .
\]
The normal equations decouple, so $\hat{\bm\beta}=(X_A\T X_A)^{-1}X_A\T\bm y_A$, $\hat{\bm\gamma}=(X_B\T X_B)^{-1}X_B\T\bm y_B$ (separate OLS fits)
and $\text{RSS}=\text{RSS}_A+\text{RSS}_B$, with $n-2p$ residual degrees of freedom. The constrained model is the pooled regression
$\bm y=X\bm\beta+\bm\varepsilon$ with residual sum $\text{RSS}_0$. By \cref{ch05:thm-F} ($q=p$ constraints $\bm\gamma-\bm\beta=\bm 0$, $2p$ parameters)
\begin{equation}\label{ch05:eq-chow}
  F=\frac{(\text{RSS}_0-\text{RSS}_A-\text{RSS}_B)/p}{(\text{RSS}_A+\text{RSS}_B)/(n-2p)}\sim F_{p,\,n-2p}\quad\text{under }H_0 .
\end{equation}
In econometrics this is the \emph{Chow test} (the name is not used in the course).

\paragraph{Approach 2: parametrise the change} \ts{24}{11}{1:11:39}. Let $\bm\delta=\bm\gamma-\bm\beta$. Then
\[
  \begin{pmatrix}\bm y_A\\ \bm y_B\end{pmatrix}=\begin{pmatrix}X_A&0\\X_B&X_B\end{pmatrix}\begin{pmatrix}\bm\beta\\ \bm\delta\end{pmatrix}+\bm\varepsilon,\qquad H_0:\bm\delta=\bm 0 .
\]
The new design is the old one times $G=\bigl(\begin{smallmatrix}I&0\\-I&I\end{smallmatrix}\bigr)$ acting on the parameters, so it has the
same column space: same fit, same RSS, same $F$ (\cref{ch05:sec-standard}). The advantage is that the individual
$t$ statistics of $\hat{\bm\delta}$ tell \emph{which} coefficients changed. In R this is the model with interactions
between the regressors and a period indicator, \texttt{y $\sim$ x * ind.periodB}.

\begin{coursenote}[Erratum in \file{lec08}, section 3.1]
The $F$ statistic of Approach~1 is printed with denominator $\text{RSS}/(n-p)$. The unconstrained model has $2p$ parameters,
so it must be $\text{RSS}/(n-2p)$ as in \eqref{ch05:eq-chow}; the text right below the formula (and section 3.2) correctly
states $df_2=n-2p$.
\end{coursenote}

\begin{casestudy}[Did GE's CAPM change? (\file{lec08}, section 4)]
\textbf{Choosing the split.} Fit the CAPM regression of \cref{ch05:ex-gebeta} to 2017--2023 and plot the cumulative
sum of the residuals standardised to equal variance, $\hat\varepsilon_i/\sqrt{1-H_{ii}}$ (\cref{ch05:sec-diag}). Under the model
this should wander like a driftless random walk; instead it falls steadily to a minimum at 2018-12-27 and then
rises \ts{24}{11}{1:12:44}. Period A: up to 2018-12-27 (500 days); period B: after.

\textbf{Fit with interactions} (\texttt{Ra.RF $\sim$ Mkt.RF * ind.periodB}, $n=1676$):

\smallskip
\centering\small
\begin{tabular}{@{}lrrr@{}}
\toprule
 & Estimate & Std.\ error & $t$\\
\midrule
intercept ($\alpha_A$) & $-0.00288$ & $0.00092$ & $-3.13$\\
\texttt{Mkt.RF} ($\beta_A$) & $0.774$ & $0.110$ & $7.02$\\
\texttt{ind.periodB} ($\delta_\alpha$) & $0.00334$ & $0.00110$ & $3.04$\\
\texttt{Mkt.RF:ind.periodB} ($\delta_\beta$) & $0.354$ & $0.118$ & $2.99$\\
\bottomrule
\end{tabular}

\smallskip
\raggedright
So $\hat\alpha_B=0.00046$ and $\hat\beta_B=1.128$: the alpha turns from negative to (slightly) positive and the beta rises
from $0.77$ to $1.13$ \ts{24}{11}{1:13:44}. The $F$ test of $\delta_\alpha=\delta_\beta=0$ (\texttt{anova(lmfit0, lmfit0a)} or
\texttt{linearHypothesis}): RSS drops from $0.71819$ to $0.71020$,
$F=\frac{0.0079821/2}{0.71020/1672}=9.40$ on $(2,1672)$ degrees of freedom, $p=8.8\times10^{-5}$.

\textbf{Lesson.} CAPM parameters are not constant over long periods; practitioners estimate them on
\emph{rolling windows}, which raises the question of the window length \ts{24}{11}{1:13:44}.
\end{casestudy}

\begin{coursenote}[Choosing the break from the data (not discussed in class)]
Here the break date was chosen \emph{after} looking at the cumulative residuals, i.e.\ at the most favourable point
for rejection. The $F_{2,1672}$ distribution assumes the date is fixed in advance, so the reported $p$-value is too
optimistic. With an unknown break date one uses the supremum of the $F$ statistics over candidate dates (Andrews'
sup-$F$ test) or CUSUM tests, which formalise exactly the cumulative-residual plot used above. The
conclusion for GE would very likely survive, given $p\approx10^{-4}$.
\end{coursenote}

\paragraph{Open issues} (\file{lec08}, section 5). The test assumes i.i.d.\ normal errors with the same variance in
both periods. If the variances differ, fit the two periods separately by maximum likelihood and use a
generalised likelihood-ratio test instead of $F$. If the errors are dependent (check with the Durbin--Watson,
Ljung--Box, turning-point, difference-sign or McLeod--Li tests), time-series models are needed
(\cref{ch:timeseries}). If they are non-normal, large-sample theory still makes $\hat{\bm\beta}$ approximately normal and
$F$ tests are replaced by $\chi^2$ tests.

\begin{casestudy}[Exchange-rate regimes of the Chinese yuan (2013 Case Study 2)]
\textbf{Question.} The yuan (CNY) was pegged to the US dollar until July 2005; China then announced a managed float
against an undisclosed \emph{basket} of currencies, with daily moves up to 0.3\%. What is the implicit basket, does it
change, and how fast does the yuan move against it?

\textbf{Data and method.} \file{fm_casestudy_fx_1.r} downloads from FRED the daily dollar exchange rates listed in
\file{fred_fxrates.txt} (23 currencies plus trade-weighted dollar indices, 1999--2013).
\file{fm_casestudy_2_rcode.r} re-expresses eight of them in Swiss francs (a stable numeraire outside the basket), so
that the dollar itself becomes a regressor, and regresses daily log changes of CNY/SFR on those of USD, JPY, EUR,
GBP, KRW, MYR and THB against the franc (the Frankel--Wei regression).

\textbf{Results.}
\begin{itemize}
  \item 2001 -- June 2005: coefficient on USD $1.000$ (se $0.0005$, $t=1839$), all others insignificant,
        $R^2=0.9999$: a perfect dollar peg.
  \item July--December 2005: USD $0.19$ (not significant), Korean won $0.18$ ($t=4.9$), Malaysian ringgit $0.75$
        ($t=5.1$), $R^2=0.949$: the peg is gone.
  \item 2006--2013, year by year: USD weight $0.87$--$0.97$ ($t$ from 25 to 97), with small, changing, sometimes
        significant weights on MYR, KRW, EUR, GBP or THB; $R^2$ between $0.93$ and $0.998$.
  \item The intercept is a daily drift against the implied basket; annualised as $(1+\hat\alpha)^{252}-1$ it is negative in
        almost every year ($-2.9\%$ in 2006, $-6.1\%$ in 2007, $-7.3\%$ in 2008, $+0.2\%$ in 2009, then $-0.5\%$ to $-4.8\%$):
        fewer yuan per franc than the basket implies, i.e.\ a steady managed \emph{appreciation} of the yuan.
  \item Diagnostics for 2012: no heteroscedasticity, but residual tails heavier than normal.
\end{itemize}
\textbf{Lesson.} A regression coefficient can identify an exchange-rate regime, and fitting it period by period reveals
regime changes; the structural break of July 2005 is visible at a glance. \file{fm_casestudy_fx_1.r} also loads the R
package \texttt{fxregime}, which implements this regression together with formal tests for changes of regime.
\end{casestudy}

\begin{exercise}[not from the course: the Chow test]
\leavevmode
\begin{enumerate}[(a)]
  \item Show that the residual sum of squares of the block-diagonal design $\bigl(\begin{smallmatrix}X_A&0\\0&X_B\end{smallmatrix}\bigr)$ is
      $\text{RSS}_A+\text{RSS}_B$.
  \item Show that $\bigl(\begin{smallmatrix}X_A&0\\X_B&X_B\end{smallmatrix}\bigr)$ has the same column space, and relate the two
      parameter vectors.
  \item For GE, compute $\hat\alpha_B$, $\hat\beta_B$ and the $F$ statistic from the numbers in the case study of
      \cref{ch05:sec-chow}. Why can you not obtain the standard error of $\hat\beta_B$ from that table alone?
\end{enumerate}
\end{exercise}

% ------------------------------------------------------------------
\section{Regression diagnostics}\label{ch05:sec-diag}
Step 4 of the modelling loop. By \cref{ch05:thm-normal}(ii) the residuals are correlated and have unequal variances,
$\Var(\hat\varepsilon_i)=\sigma^2(1-H_{ii})$, even when the errors are i.i.d.\ \ts{24}{8}{1:01:33}. Diagnostics correct for this.

\paragraph{Leverage.} $H_{ii}$ is the \emph{leverage} of case $i$. Since $\hat y_i=\sum_jH_{ij}y_j$, $\partial\hat y_i/\partial y_i=H_{ii}$: the
fraction of its own fitted value that a case dictates (2013 PS3). From $H=H^2=H\T$,
$H_{ii}=\sum_jH_{ij}^2\ge H_{ii}^2$, so $0\le H_{ii}\le1$; with an intercept $H_{ii}\ge1/n$, and $\sum_iH_{ii}=p$, so the average is
$p/n$ \ts{24}{8}{1:05:44}. If $H_{ii}=1$ the fitted line passes through the point: that case alone determines its
fit \ts{24}{8}{1:04:44}. With an intercept,
\[
  H_{ii}=\frac1n+(\bm x_i-\bar{\bm x})\T(\mathcal X\T\mathcal X)^{-1}(\bm x_i-\bar{\bm x}),
\]
where $\mathcal X$ is the matrix of centred covariates (2013 PS3, \cref{ch05:ex-ps3-1de}): leverage grows with the
Mahalanobis distance of $\bm x_i$ from the centroid of the covariates (dividing by $n-1$ turns $\mathcal X\T\mathcal X$ into the sample
covariance matrix). High leverage is a property of the $x$'s only. R's \texttt{influence.measures()} flags cases with
$H_{ii}>3p/n$.

\paragraph{Residuals.} The \emph{standardised} residual (R: \texttt{rstandard}) is
$r_i=\hat\varepsilon_i/(\hat\sigma\sqrt{1-H_{ii}})$; the \emph{studentised} (deleted) residual (R: \texttt{rstudent}, or
\texttt{studres} in \texttt{MASS}) uses the estimate $\hat\sigma_{(i)}$ computed without case $i$,
$t_i=\hat\varepsilon_i/(\hat\sigma_{(i)}\sqrt{1-H_{ii}})$, and has exactly a $t_{n-p-1}$ distribution under the normal model, which makes
outliers easy to judge \ts{24}{8}{1:02:36}.

\begin{proposition}[Case deletion]\label{ch05:prop-deletion}
Let $\hat{\bm\beta}_{(i)}$, $\hat{\bm y}_{(i)}=X\hat{\bm\beta}_{(i)}$ and $\hat\sigma^2_{(i)}$ be computed with case $i$ deleted. Then
\[
  \hat{\bm\beta}_{(i)}=\hat{\bm\beta}-\frac{(X\T X)^{-1}\bm x_i\hat\varepsilon_i}{1-H_{ii}},\qquad
  \hat\sigma^2_{(i)}=\hat\sigma^2+\frac{1}{n-p-1}\Bigl(\hat\sigma^2-\frac{\hat\varepsilon_i^2}{1-H_{ii}}\Bigr),
\]
and \emph{Cook's distance} is
\[
  CD_i=\frac{\|\hat{\bm y}-\hat{\bm y}_{(i)}\|^2}{p\,\hat\sigma^2}=\frac{\hat\varepsilon_i^2}{p\,\hat\sigma^2}\,\frac{H_{ii}}{(1-H_{ii})^2}
  =\frac{r_i^2}{p}\,\frac{H_{ii}}{1-H_{ii}} .
\]
\end{proposition}
The proofs, via the Sherman--Morrison--Woodbury formula (deleting a row is a rank-one update of $X\T X$), are 2013 PS3
Problem~2 (\cref{ch05:ex-ps3-2a,ch05:ex-ps3-2bd}). Nothing needs to be refitted: all leave-one-out quantities come
from one fit. The last form of $CD_i$ says \emph{influence = outlyingness $\times$ leverage}: a large residual at a
low-leverage point, or a high-leverage point on the line, is harmless; both together move the fit. Since
$\|\hat{\bm y}-\hat{\bm y}_{(i)}\|^2=(\hat{\bm\beta}-\hat{\bm\beta}_{(i)})\T X\T X(\hat{\bm\beta}-\hat{\bm\beta}_{(i)})$, $CD_i$ measures how far $\hat{\bm\beta}$ moves in the metric of
its own confidence ellipsoid: values comparable with the median of $F_{p,n-p}$ (about $1$) are large
\ts{24}{8}{1:06:44}.

\paragraph{Graphical tools} \ts{24}{8}{1:08:51}. In R (\texttt{stats} and
\texttt{car}):
\begin{center}
\small
\begin{tabularx}{\linewidth}{@{}l L@{}}
\toprule
\texttt{plot(fit)} & residuals vs fitted (flat band), normal Q--Q,
scale--location ($\sqrt{|r_i|}$ vs fitted), residuals vs leverage
with Cook contours\\
\texttt{influence.measures()} & leave-one-out stats: DFBETAS,
DFFITS, COVRATIO, Cook's distance, hat values\\
\texttt{qqPlot()} & Q--Q plot of studentised residuals against
$t$ quantiles, with confidence band \ts{24}{8}{1:03:44}\\
\texttt{influencePlot()} & studentised residuals vs hat values,
bubble size $\propto$ Cook's distance\\
\texttt{residualPlots()} & residuals vs each regressor; curvature
test = $t$ test of added squared term; Tukey test = same with
$\hat y^2$ \ts{24}{8}{1:09:53}\\
\texttt{leveragePlots()} & added-variable plots: effect of each
regressor after adjusting for others\\
\texttt{durbinWatsonTest()} & autocorrelation of residuals
(time-ordered data)\\
\bottomrule
\end{tabularx}
\end{center}

\begin{casestudy}[Prostate cancer data (\file{lec06_1}, slides 21--33)]
\textbf{Data.} Not financial, deliberately: the \texttt{prostate} data of the \texttt{faraway} package (used in Hastie,
Tibshirani and Friedman's \emph{Elements of Statistical Learning}), 97 patients, response \texttt{lpsa} (log PSA), eight
covariates (log cancer volume \texttt{lcavol}, log prostate weight, age, \dots, seminal vesicle invasion \texttt{svi}).
A random 75\% training sample ($n=72$, \texttt{set.seed(1)}) is used. The experience transfers: only the variable
names change from one regression problem to another \ts{24}{8}{46:42}.

\textbf{Fit} ($p=9$ with intercept): $R^2=0.630$, $R^2_{\text{adj}}=0.583$, $\hat\sigma=0.724$ on 63 degrees of freedom,
$F=13.4$. Significant: \texttt{lcavol} ($t=5.21$) and \texttt{svi} ($t=2.95$); \texttt{lweight} borderline ($p=0.07$). The pairs
plot already shows the strong \texttt{lpsa}--\texttt{lcavol} relation \ts{24}{8}{48:48}. Standardising the covariates
leaves every $t$ value unchanged (\cref{ch05:sec-standard}) \ts{24}{8}{57:24}.

\textbf{Diagnostics.} Studentised residuals: unimodal, symmetric, within the Q--Q band except case 39 ($-3.06$). Hat
values: cases 32 ($0.41$) and 74 ($0.33$) exceed $3p/n=0.375$ or come close, against an average $p/n=0.125$. Cook's
distances of the flagged cases are only $0.07$--$0.13$. Curvature tests: none significant (smallest $p=0.10$ for
\texttt{svi}; Tukey test $p=0.69$) \ts{24}{8}{1:09:53}.

\textbf{Lesson.} A satisfactory model: no curvature, no dominant case; the scale--location plot hints at mild
non-constancy of the variance.
\end{casestudy}

\begin{exercise}[PS3 (2013), Problem 1(d)--(e): invariance and leverage]\label{ch05:ex-ps3-1de}
(d) Prove that $H$ is unchanged if $X$ is replaced by $X'=XG$ for any invertible $p\times p$ matrix $G$.
(e) Let $X$ be $n\times(p+1)$ with a first column of ones and $p$ covariates, and let $G$ be the identity except for the first
row, which is $(1,-\bar x_1,\dots,-\bar x_p)$, with $\bar x_j=\frac1n\sum_ix_{i,j}$. By (d) the models with $X$ and $X'=XG$ have the
same fitted values and residuals.
\begin{itemize}
  \item If $\bm\beta=(\beta_0,\dots,\beta_p)\T$ is the parameter for $X$, show that $\bm\beta'=G^{-1}\bm\beta$ is the parameter for $X'$; find
        $G^{-1}$ and the elements of $\bm\beta'$ explicitly.
  \item Show that $X'^{\mathsf T}X'=\bigl(\begin{smallmatrix}n&\bm 0_p\T\\ \bm 0_p&\mathcal X\T\mathcal X\end{smallmatrix}\bigr)$, where $\mathcal X$ has entries $x_{i,j}-\bar x_j$.
  \item Prove $H_{ij}=\frac1n+(\bm x_i-\bar{\bm x})\T[\mathcal X\T\mathcal X]^{-1}(\bm x_j-\bar{\bm x})$, where $\bm x_i=(x_{i,1},\dots,x_{i,p})\T$. Hence the
        leverage $H_{ii}$ increases with the (squared) Mahalanobis distance between $\bm x_i$ and $\bar{\bm x}$.
\end{itemize}
\end{exercise}

% ------------------------------------------------------------------
\section{Shrinkage and principal-components regression}\label{ch05:sec-shrink}
When regressors are many or strongly correlated (the four index ETFs below have pairwise correlations
$0.77$--$0.95$), $X\T X$ is ill-conditioned and $\Cov(\hat{\bm\beta})=\sigma^2(X\T X)^{-1}$ is huge in some directions. Three
remedies, all expressed through the SVD, were presented in L11.

\paragraph{Standardisation first} \ts{24}{11}{23:47}. Centre $\bm y$ and the columns of $X$ and scale the columns to unit
variance; the intercept is then estimated by $\bar y$. Otherwise a penalty on $\|\bm\beta\|$ would depend on the units:
measuring a covariate in cents instead of dollars would make its coefficient 100 times smaller and almost free
\ts{24}{11}{24:49}. Below $X$ is the standardised $n\times p$ matrix (no intercept column) with thin SVD
$X=UDV\T$, $U\T U=I_p$, $V$ orthogonal, $D=\diag(d_1\ge\dots\ge d_p)$ (\cref{thm:svd}), and $c_j=\bm u_j\T\bm y$.

\begin{coursenote}[Erratum and notation in the ridge slides]
Slide 41 of \file{lec06_1} rescales the columns by ``$S^{-1/2}X$'' with $S=\frac1n\diag(X\T X)$: since $S$ is $p\times p$ this must be
$XS^{-1/2}$ (multiplication on the right scales columns). On the same slide $H^*=I_n-\frac1n\bm 1\bm 1\T$ is the
\emph{centring} matrix, not the hat matrix.
\end{coursenote}

\subsection{Ridge regression}
\begin{definition}[Ridge]
For a complexity parameter $\lambda\ge0$ \ts{24}{11}{22:46},
\[
  \hat{\bm\beta}^{\text{ridge}}=\arg\min_{\bm\beta}\Bigl(\|\bm y-X\bm\beta\|^2+\lambda\|\bm\beta\|^2\Bigr)=(X\T X+\lambda I_p)^{-1}X\T\bm y .
\]
\end{definition}
The gradient condition is $-2X\T(\bm y-X\bm\beta)+2\lambda\bm\beta=\bm 0$; $X\T X+\lambda I$ is invertible for every $\lambda>0$ even when $X$ is
rank-deficient. The estimate is linear in $\bm y$ and equals OLS for $\lambda=0$. With the SVD
\ts{24}{11}{28:04}:
\[
  \hat{\bm\beta}^{\text{ridge}}=V(D^2+\lambda I)^{-1}DU\T\bm y,\qquad
  \hat{\bm y}^{\text{ridge}}=UD(D^2+\lambda I)^{-1}DU\T\bm y=\sum_{j=1}^p\frac{d_j^2}{d_j^2+\lambda}\,c_j\,\bm u_j ,
\]
whereas $\hat{\bm y}^{\text{LS}}=UU\T\bm y=\sum_jc_j\bm u_j$. Ridge shrinks the component of the fit along $\bm u_j$ by the factor
$d_j^2/(d_j^2+\lambda)$: little along the directions in which the regressors vary a lot (large $d_j$, the leading principal
components), much along the directions of little variation \ts{24}{11}{29:04}. It is biased,
$\E\hat{\bm\beta}^{\text{ridge}}=(X\T X+\lambda I)^{-1}X\T X\bm\beta$, but has smaller variance; $\lambda$ is chosen by cross-validation
(\texttt{cv.glmnet} in the case study).

\begin{proposition}[Ridge as a Bayes estimate]\label{ch05:prop-ridgebayes}
In the normal linear model with prior $\bm\beta\sim N_p(\bm 0,\tau^2I_p)$, the posterior mode (which, the posterior being Gaussian,
is also its mean) is $\hat{\bm\beta}^{\text{ridge}}$ with $\lambda=\sigma^2/\tau^2$ \ts{24}{11}{25:52}.
\end{proposition}
\begin{proof}
The posterior density is proportional to likelihood $\times$ prior,
\[\exp\bigl(-\|\bm y-X\bm\beta\|^2/(2\sigma^2)\bigr)\exp\bigl(-\|\bm\beta\|^2/(2\tau^2)\bigr);\]
maximising it means minimising $\|\bm y-X\bm\beta\|^2+(\sigma^2/\tau^2)\|\bm\beta\|^2$.
\end{proof}
The isotropic prior says that a priori no direction of $\R^p$ is preferred \ts{24}{11}{26:58}; a strong prior (small
$\tau$, large $\lambda$) pulls all coefficients towards zero.

\begin{exercise}[not from the course: ridge as augmented least squares]
\leavevmode
\begin{enumerate}[(a)]
  \item Show that $\hat{\bm\beta}^{\text{ridge}}$ is the OLS estimate for the augmented data $\tilde{\bm y}=(\bm y,\bm 0_p)$,
      $\tilde X=\bigl(\begin{smallmatrix}X\\ \sqrt\lambda I_p\end{smallmatrix}\bigr)$.
  \item Using the SVD, show that $\|\hat{\bm\beta}^{\text{ridge}}\|$ decreases in $\lambda$ and compute
      $\Cov(\hat{\bm\beta}^{\text{ridge}})$ under Gauss--Markov.
  \item Show that ridge has smaller variance than OLS in every direction:
      $\Cov(\hat{\bm\beta})-\Cov(\hat{\bm\beta}^{\text{ridge}})\succeq0$.
\end{enumerate}
\end{exercise}

\subsection{Principal-components regression}
The principal components of the (centred) regressors are $\bm z_j=X\bm v_j$, where $\bm v_j$ are the eigenvectors of the
sample covariance $S=\frac1nX\T X=\frac1nVD^2V\T$ \ts{24}{11}{30:05}. They are orthogonal, $\bm z_j=d_j\bm u_j$, with sample
variance $d_j^2/n$, and $\bm z_1$ is the unit-norm combination of the columns of $X$ with the largest variance
(\cref{ch:pca}). \emph{PC regression} fixes $m\le p$ and regresses $\bm y$ on $Z_m=[\bm z_1\cdots\bm z_m]=U_mD_m$
\ts{24}{11}{31:05}:
\[
  \hat{\bm\gamma}=(Z_m\T Z_m)^{-1}Z_m\T\bm y,\quad \hat\gamma_j=\frac{\bm z_j\T\bm y}{d_j^2},\qquad
  \hat{\bm\beta}^{\text{pc}}=V_m\hat{\bm\gamma},\qquad \hat{\bm y}^{\text{pc}}=U_mU_m\T\bm y=\sum_{j=1}^mc_j\bm u_j .
\]
Orthogonality makes each $\hat\gamma_j$ a separate simple regression, and
\[\Cov(\hat{\bm\beta}^{\text{pc}})=\sigma^2V_mD_m^{-2}V_m\T\preceq\sigma^2VD^{-2}V\T=\Cov(\hat{\bm\beta}):\]
variance is reduced by dropping the noisiest
directions \ts{24}{11}{32:13}.

\paragraph{Comparison} \ts{24}{11}{33:19}. The three fits use the same building blocks $c_j\bm u_j$:
\[
  \hat{\bm y}^{\text{LS}}=\sum_{j=1}^pc_j\bm u_j,\qquad
  \hat{\bm y}^{\text{ridge}}=\sum_{j=1}^p\frac{d_j^2}{d_j^2+\lambda}c_j\bm u_j,\qquad
  \hat{\bm y}^{\text{pc}}=\sum_{j=1}^mc_j\bm u_j .
\]
Least squares keeps all, ridge shrinks smoothly, PC regression keeps the first $m$ and discards the rest.

\begin{intuition}[Spectral filters]
This is the classic regularisation of an ill-posed inverse problem: ridge is Tikhonov regularisation (a smooth
filter $d^2/(d^2+\lambda)$ on the singular values), PC regression is the truncated SVD (a sharp cut-off). Both trade a
little bias for a large reduction of the noise amplified by small singular values.
\end{intuition}

A caveat, visible in the ETF case study \ts{24}{11}{46:50}: the principal components are ranked by how much the
\emph{regressors} vary, not by how much they explain $\bm y$. A low-variance component can be an important predictor,
so keeping ``the first $m$'' may throw away the useful direction.

\subsection{The lasso}
\begin{definition}[Lasso]
\[
  \hat{\bm\beta}^{\text{lasso}}=\arg\min_{\bm\beta}\Bigl(\|\bm y-X\bm\beta\|^2+\lambda\sum_{j=1}^p|\beta_j|\Bigr),
\]
equivalently: minimise $\|\bm y-X\bm\beta\|^2$ subject to $\|\bm\beta\|_1\le t$ \ts{24}{11}{33:19}. The constraint binds only for
$t<t_0=\|\hat{\bm\beta}^{\text{LS}}\|_1$; the \emph{regularisation path} is $\hat{\bm\beta}^{\text{lasso}}$ as a function of $\lambda\ge0$ (or of
$t\in[0,t_0]$).
\end{definition}
The $L_1$ ball is a diamond (a cross-polytope in $\R^p$) with corners on the axes. The elliptical contours of
$\|\bm y-X\bm\beta\|^2$, centred at $\hat{\bm\beta}^{\text{LS}}$, typically first touch it at a corner, where some coefficients are exactly
zero \ts{24}{11}{35:28} (\cref{ch05:fig-ridgelasso}). The lasso therefore does \emph{variable selection}: as $\lambda$ grows,
variables leave the model one after the other \ts{24}{11}{36:31}. There is no closed form; the path is computed by
coordinate descent (\texttt{glmnet}).

\begin{figure}[ht]
\centering
\begin{tikzpicture}[scale=0.95]
  \begin{scope}
    \draw[->,black!60] (-1.6,0) -- (2.8,0) node[right,black]{\small $\beta_1$};
    \draw[->,black!60] (0,-1.5) -- (0,4.3) node[above,black]{\small $\beta_2$};
    \fill[intcol!25] (0,0) circle (1);
    \draw[intcol!80!black,thick] (0,0) circle (1);
    \foreach \s in {2.093,1.35,0.6}
      \draw[thmcol,rotate around={-50:(0.6,2.2)}] (0.6,2.2) ellipse ({\s} and {0.5*\s});
    \fill[thmcol] (0.6,2.2) circle (1.5pt) node[right]{\small $\hat{\bm\beta}^{\text{LS}}$};
    \fill[black] (0.562,0.827) circle (1.8pt);
    \node[anchor=north west,font=\small] at (0.62,0.78) {$\hat{\bm\beta}^{\text{ridge}}$};
    \node[font=\small] at (0.6,-1.25) {ridge: $\beta_1^2+\beta_2^2\le t^2$};
  \end{scope}
  \begin{scope}[xshift=6.3cm]
    \draw[->,black!60] (-1.6,0) -- (2.8,0) node[right,black]{\small $\beta_1$};
    \draw[->,black!60] (0,-1.5) -- (0,4.3) node[above,black]{\small $\beta_2$};
    \fill[defcol!25] (1,0) -- (0,1) -- (-1,0) -- (0,-1) -- cycle;
    \draw[defcol,thick] (1,0) -- (0,1) -- (-1,0) -- (0,-1) -- cycle;
    \foreach \s in {2.519,1.6,0.7}
      \draw[thmcol,rotate around={-50:(0.6,2.2)}] (0.6,2.2) ellipse ({\s} and {0.5*\s});
    \fill[thmcol] (0.6,2.2) circle (1.5pt) node[right]{\small $\hat{\bm\beta}^{\text{LS}}$};
    \fill[black] (0,1) circle (1.8pt);
    \node[anchor=east,font=\small] at (-0.08,1.18) {$\hat{\bm\beta}^{\text{lasso}}$};
    \node[font=\small] at (0.6,-1.25) {lasso: $|\beta_1|+|\beta_2|\le t$};
  \end{scope}
\end{tikzpicture}
\caption{Constrained forms of ridge and lasso (same least-squares contours, $t=1$). The solution is where the smallest
contour touches the constraint region: a generic point of the circle for ridge, the corner $\beta_1=0$ of the diamond for
the lasso \ts{24}{11}{34:22}. (Contours computed for a quadratic form with axis ratio $2{:}1$.)}\label{ch05:fig-ridgelasso}
\end{figure}

\begin{remark}[Orthonormal design; not discussed in class]
If $X\T X=I_p$, then $\|\bm y-X\bm\beta\|^2=\|\bm y-X\hat{\bm\beta}\|^2+\|\bm\beta-\hat{\bm\beta}\|^2$ with $\hat{\bm\beta}=X\T\bm y$, and both problems decouple:
\[
  \hat\beta^{\text{ridge}}_j=\frac{\hat\beta_j}{1+\lambda},\qquad
  \hat\beta^{\text{lasso}}_j=\operatorname{sign}(\hat\beta_j)\bigl(|\hat\beta_j|-\lambda/2\bigr)_+ .
\]
Ridge rescales every coefficient; the lasso subtracts a constant and sets the small ones to exactly zero
(``soft thresholding'').
\end{remark}

Regularised models reappear in the 2013 lecture on regularised pricing and risk models (\cref{ch:regularized}), in
machine learning (\cref{ch:ml}), and PCA in \cref{ch:pca}.

% ------------------------------------------------------------------
\section{Case studies}\label{ch05:sec-cases}

\subsection{Sector ETFs regressed on index ETFs}
\begin{casestudy}[ETF case study (\file{lec06_4} = \file{ETF_Casestudy.pdf}, sector XLP)]
\textbf{Data.} Adjusted daily prices from Yahoo (2000-01-03 to 2024-09-19) of the 13 SPDR ETFs listed in
\file{spyders_symbols_labels.csv}: four index ETFs (SPY: S\&P~500, MDY: mid-caps, QQQ: Nasdaq-100, DIA: Dow Jones) and
nine sector ETFs (XLB materials, XLV health care, XLP consumer staples, XLY consumer discretionary, XLE energy, XLF
financials, XLI industrials, XLK technology, XLU utilities). Converted to weekly log returns (1289 weeks; saved in
\file{casestudy_1_0_etfs.RData}); the analysis uses 2013--2024 ($n=612$ weeks). Response: XLP (annualised mean 10.0\%,
volatility 13.7\%, against 13.7\% and 16.1\% for SPY).

\textbf{OLS} \ts{24}{11}{42:42}: $\hat y=0.0002+1.007\,\text{SPY}-0.368\,\text{MDY}-0.314\,\text{QQQ}+0.400\,\text{DIA}$, all
slopes with $|t|>4.3$; $R^2=0.642$, $\hat\sigma=0.0114$ (weekly), $F=271.9$. The mixed signs are a symptom of the collinear
regressors (index correlations $0.77$--$0.95$): XLP loads positively on the market but less on its growth and mid-cap
parts. Diagnostics: influential weeks 2016-11-11 (studentised residual $-4.8$), 2020-03-06 ($+4.4$), 2022-05-20
($-5.8$); the COVID week 2020-03-20 has hat value $0.135$ against $p/n=0.008$. The curvature tests are all significant
at 5\% ($p\approx0.02$--$0.04$): some non-linearity in the market response.

\textbf{$z$-scores.} Regressing on standardised index returns reproduces every $t$ value and $R^2$; only the scale of the
coefficients and the intercept change.

\textbf{PCA of the regressors} \ts{24}{11}{43:42}: the four standardised index returns have PC variances in proportions
$90.9\%$, $6.5\%$, $2.3\%$, $0.3\%$ (cumulative $97.4\%$ with two, $99.7\%$ with three). PC1 has equal loadings
($0.48$--$0.52$, the ``market''), PC2 contrasts QQQ ($+0.79$) with MDY and DIA.

\textbf{PC regressions} \ts{24}{11}{44:48}: simple regressions of XLP on each PC have $R^2=0.545$, $0.006$ ($p=0.052$),
$0.079$ and $0.012$. Their sum is $0.642$, the $R^2$ of the full regression, because the PCs are orthogonal
(\cref{ch05:ex-r2orth}). The third component, with only 2.3\% of the regressor variance, explains 13 times more of
XLP than the second: dropping it (PC123 would keep it, but a ``first two'' rule would not) changes the answer. Using PCs 1--3
gives coefficients very different from OLS; using the PCs with $|t|>\sqrt2$ (all four here) reproduces OLS exactly
\ts{24}{11}{51:03}.

\textbf{Ridge and lasso} (\texttt{glmnet}, $\lambda$ by 10-fold cross-validation over $10^{-5}$--$10^{3}$)
\ts{24}{11}{47:55}: coefficient traces shrink to zero, not monotonically for ridge; along the lasso path the third
variable (QQQ) leaves first, then MDY, then DIA \ts{24}{11}{50:01}. With the CV-optimal $\lambda$ both fits are
almost identical to OLS ($R^2=0.6417$ ridge, $0.6417$ lasso vs $0.6418$): with $n=612$ and 4 regressors there is little
to regularise \ts{24}{11}{51:03}.

\textbf{XLB version} (\file{ETF_Casestudy_XLB.pdf}, materials): $R^2=0.805$, loadings SPY $0.85$, MDY $0.50$, QQQ
$-0.38$, DIA not significant ($t=0.43$); PC1 alone gives $R^2=0.73$. The week 2020-03-20 is now very influential (Cook's
distance $0.99$), and all curvature tests are highly significant ($p<0.001$).

\textbf{Remark on the code.} Cumulative returns are plotted as $\prod_t(1+r_t)$ although $r_t$ are log returns; the correct
growth factor is $\exp(\sum_tr_t)$. The difference (about $-r_t^2/2$ per week) only affects the plots.
\end{casestudy}

\begin{finance}[Why regress a sector on the indices?]
The fitted coefficients are a \emph{hedge}: holding XLP and shorting $1.007$ SPY, $-0.368$ MDY, \dots\ leaves (in
sample) only the residual, i.e.\ the sector-specific part of the return, with most of the market risk removed
\ts{24}{11}{40:39}. The same regression, run with a hedge fund's returns as the response and liquid instruments as
regressors, is the basis of \emph{hedge-fund replication}: if most of a fund's return is explained by tradable factors,
those can be bought cheaply \ts{24}{11}{41:39}. (Regressing an index on its sectors would instead give $R^2\approx1$, as
noted by a student \ts{24}{11}{39:39}.)
\end{finance}

\subsection{CAPM regressions}\label{ch05:capm}
The CAPM (Sharpe 1964, Lintner 1965; theory in \cref{ch:portfolio}) states that in an efficient market in which all
investors hold mean--variance efficient portfolios and can borrow and lend at the risk-free rate $r_f$,
\[
  \E[R_j]=r_f+\beta_j\bigl(\E[R_M]-r_f\bigr),\qquad \beta_j=\frac{\Cov(R_j,R_M)}{\Var(R_M)},
\]
where $R_M$ is the return of the market portfolio \ts{24}{11}{54:12}. With excess returns $R^*_{j,t}=R_{j,t}-r_{f,t}$ and
$R^*_{M,t}$ observed over time, fit
\[
  R^*_{j,t}=\alpha_j+\beta_jR^*_{M,t}+\varepsilon_{j,t},\qquad \varepsilon_{j,t}\ \text{white noise},
\]
the alpha--beta regression of \cref{def:alphabeta} \ts{24}{11}{55:18}. Taking expectations, the CAPM holds iff
$\alpha_j=0$ \ts{24}{11}{56:27}. Two natural tests: $\alpha_j=0$ (is the asset efficiently priced?) and $\beta_j=1$ (does it carry
more or less systematic risk than the market?) \ts{24}{11}{1:06:14}.

\begin{casestudy}[CAPM and macro factors for GE and XOM, 2000--2013 (2013 Case Study 1)]
\textbf{Data.} \file{fm_casestudy_1_0.r} downloads daily prices of BAC, GE, JDSU, XOM and the S\&P~500 (Yahoo) and from FRED
the constant-maturity Treasury yields (3 months, 1, 5, 10 years), Moody's Aaa and Baa yields and the WTI oil price, for
2000-01-03 to 2013-05-31 (3373 days). \file{fm_casestudy_1_rcode.r} computes daily log returns and the daily risk-free
return $\log\bigl(1+r_{3M}\,\Delta t/360\bigr)$ with $\Delta t$ the number of calendar days since the previous close (money-market
convention).

\textbf{GE:} $\hat\alpha=-0.00013$ ($t=-0.56$), $\hat\beta=1.184$ ($t=66.6$), $R^2=0.568$: consistent with CAPM, with more market risk
than average. Influence: 129 of 3372 days exceed the leverage threshold; the plot of hat values over time shows they are
concentrated in the 2008 crisis, when market moves were extreme.

\textbf{Adding oil} (daily log change of WTI): for GE the oil coefficient is $-0.034$ ($t=-3.6$): GE falls when oil rises. For
XOM, simple CAPM $\hat\beta=0.830$, $R^2=0.451$ (lower than GE: an oil company tracks the market less); with oil,
$\hat\beta=0.782$ and oil coefficient $+0.132$ ($t=16.1$), $R^2$ up to $0.491$. The high-leverage days are those with extreme
market or oil moves, far from the centre of the $(R^*_M,\text{oil})$ scatter (the Mahalanobis interpretation of 2013
PS3).

\textbf{Lesson.} The market factor explains half of a large stock's daily variance; a second factor can be highly
significant while adding little $R^2$.
\end{casestudy}

\begin{casestudy}[CAPM for GE and for the S\&P~500 stocks (2024 L11)]
\textbf{GE 2017--2023} (\file{lec08}; excess return \texttt{Ra.RF} regressed on the market excess return
\texttt{Mkt.RF} of the Fama--French data): see \cref{ch05:ex-gebeta} and the regime test of \cref{ch05:sec-chow}. The residuals
are clearly non-normal \ts{24}{11}{1:00:53}; see the box on heavy tails in \cref{ch05:sec-mle}.

\textbf{All S\&P~500 stocks} (note shown in class, not published) \ts{24}{11}{1:14:45}: the same regression for the
380--400 constituents with a long enough history, grouped by sector. $R^2$ ranges over $(0,1)$, typically above $0.2$.
Alphas vary by sector (positive for construction and computer technology, negative for conglomerates)
\ts{24}{11}{1:15:47}, but the $p$-values of $\alpha_j=0$ are above $0.05$ for all but a handful of stocks, roughly the 5\%
expected by chance: the CAPM is not rejected stock by stock \ts{24}{11}{1:16:47}. Higher betas tend to go with higher
alphas \ts{24}{11}{1:17:53}. Betas are strongly sector-dependent: consumer staples lowest (the ten lowest betas, $0.37$
down to $0.18$, are mostly staples), technology and consumer discretionary highest \ts{24}{11}{1:18:59}. The largest
significant alpha, ENPH (energy), is $0.003$ per day, i.e.\ $(1.003)^{252}-1\approx113\%$ per year
\ts{24}{11}{1:20:01}.

\textbf{Next step.} Multi-factor models with the Fama--French factors, already downloaded with the data
\ts{24}{11}{1:22:11} (\cref{ch:pca}).
\end{casestudy}

\begin{finance}[Using alpha and beta]
Beta is the hedge ratio against the index and the measure of systematic risk that the CAPM says is rewarded;
alpha is the part of the return not explained by market exposure (\cref{def:alphabeta}). In practice betas are estimated
on rolling windows (\cref{ch05:sec-chow}) and often shrunk towards $1$, e.g.\ the widely quoted ``adjusted beta''
$\frac23\hat\beta+\frac13$ (not discussed in class): the same bias--variance trade as ridge regression, with the prior
centred at the market beta.
\end{finance}

% ------------------------------------------------------------------
\section{Differences between the 2013 and 2024 lectures}
\begin{coursenote}[2013 L6 vs 2024 L6, L8, L11]
The core (OLS, hat matrix, Gauss--Markov, GLS, normal distribution theory, ML, M-estimators) is the same material in
both years, with nearly identical slides. \textbf{2013} spends one lecture on it and adds the QR decomposition and
the full proofs of the Gauss--Markov and distribution theorems \ts{13}{6}{1:08:26}; its case studies are the CAPM with
oil (Case Study~1) and the yuan regimes (Case Study~2). \textbf{2024} spreads the topic over parts of three lectures and
adds real-data diagnostics (prostate data), hypothesis testing and regime-change tests (\file{lec08}), ridge, PC and lasso
regression, the ETF case study and the S\&P~500 CAPM study. The 2013 lecture ran out of time before M-estimation
\ts{13}{6}{1:21:11}; it is covered in the notes and in 2024 L11.
\end{coursenote}

% ------------------------------------------------------------------
\begin{keyformulas}
\begin{align*}
&\text{OLS:} && \hat{\bm\beta}=(X\T X)^{-1}X\T\bm y=R^{-1}Q\T\bm y,\\
  &&& H=X(X\T X)^{-1}X\T=QQ\T,\quad \hat{\bm\varepsilon}=(I-H)\bm y\\
&\text{Geometry:} && H^2=H=H\T,\ \tr H=p,\\
  &&& X\T\hat{\bm\varepsilon}=\bm 0,\ \text{TSS}=\text{ESS}+\text{RSS},\ R^2=\Corr(\bm y,\hat{\bm y})^2\\
&\text{Gauss--Markov:} && \E\hat{\bm\beta}=\bm\beta,\ \Cov\hat{\bm\beta}=\sigma^2(X\T X)^{-1};\\
  &&& \bm c\T\hat{\bm\beta}\ \text{is BLUE}\\
&\text{GLS:} && \hat{\bm\beta}_{\text{GLS}}=(X\T\Sigma^{-1}X)^{-1}X\T\Sigma^{-1}\bm y,\\
  &&& \Cov=\sigma^2(X\T\Sigma^{-1}X)^{-1}\\
&\text{Normal model:} && \hat{\bm\beta}\sim N_p(\bm\beta,\sigma^2(X\T X)^{-1})\\
  &&& \perp\!\!\!\perp\text{RSS}\sim\sigma^2\chi^2_{n-p},\ \hat\sigma^2=\tfrac{\text{RSS}}{n-p}\\
&t\text{ test:} && \hat t_j=(\hat\beta_j-b)/(\hat\sigma\sqrt{C_{jj}})\sim t_{n-p},\quad C_{jj}=[(X\T X)^{-1}]_{jj}\\
&F\text{ test:} && F=\tfrac{(\text{RSS}_0-\text{RSS})/q}{\text{RSS}/(n-p)}\sim F_{q,n-p}\\
&&& \text{RSS}_0-\text{RSS}=(A\hat{\bm\beta}-\bm b)\T[A(X\T X)^{-1}A\T]^{-1}(A\hat{\bm\beta}-\bm b)\\
&\text{Chow test:} && F=\tfrac{(\text{RSS}_0-\text{RSS}_A-\text{RSS}_B)/p}{(\text{RSS}_A+\text{RSS}_B)/(n-2p)}\\
  &&& \sim F_{p,n-2p}\\
&\text{ML:} && \hat{\bm\beta}_{\text{ML}}=\hat{\bm\beta},\quad \hat\sigma^2_{\text{ML}}=\text{RSS}/n\\
&\text{Diagnostics:} && \Var\hat\varepsilon_i=\sigma^2(1-H_{ii}),\ \ CD_i=\tfrac{r_i^2}{p}\tfrac{H_{ii}}{1-H_{ii}},\\
  &&& \hat{\bm\beta}_{(i)}=\hat{\bm\beta}-\tfrac{(X\T X)^{-1}\bm x_i\hat\varepsilon_i}{1-H_{ii}}\\
&\text{Shrinkage:} && \hat{\bm y}^{\text{ridge}}=\textstyle\sum_j\frac{d_j^2}{d_j^2+\lambda}c_j\bm u_j,\\
  &&& \hat{\bm y}^{\text{pc}}=\sum_{j\le m}c_j\bm u_j,\\
  &&& \lambda=\sigma^2/\tau^2\ \text{(Bayes)}\\
&\text{CAPM:} && R^*_{j,t}=\alpha_j+\beta_jR^*_{M,t}+\varepsilon_{j,t},\quad H_0:\alpha_j=0
\end{align*}
\end{keyformulas}

% ------------------------------------------------------------------
\section*{Exercises}
\addcontentsline{toc}{section}{Exercises}
Standalone exercises follow the section they practise; the connected series (2013 PS3, Problems~2 and~3, which build on one another) are kept together here.

\subsection*{From 2013 Problem Set 3}
Throughout, $\bm y=X\bm\beta+\bm\varepsilon$ with $X\in\R^{n\times p}$ of full rank $p\le n$, $\hat{\bm\beta}=(X\T X)^{-1}X\T\bm y$, $H=X(X\T X)^{-1}X\T$,
$\hat{\bm\varepsilon}=(I_n-H)\bm y$.

\begin{exercise}[PS3 (2013), Problem 2(a): Sherman--Morrison--Woodbury]\label{ch05:ex-ps3-2a}
Let $A$ be a $p\times p$ symmetric matrix of rank $p$ and $a,b$ two $q\times p$ matrices of rank $q<p$. Prove that, whenever the
inverses exist,
\[
  (A+a\T b)^{-1}=A^{-1}-A^{-1}a\T\bigl(I_q+bA^{-1}a\T\bigr)^{-1}bA^{-1}.
\]
\end{exercise}

\begin{exercise}[PS3 (2013), Problem 2(b)--(d): case-deletion influence]\label{ch05:ex-ps3-2bd}
(b) Using Sherman--Morrison--Woodbury, show that the least-squares estimate with case $i$ deleted is
$\hat{\bm\beta}_{(i)}=\hat{\bm\beta}-(X\T X)^{-1}\bm x_i\hat\varepsilon_i/(1-H_{ii})$, where $\bm x_i\T$ is the $i$-th row of $X$ and $\hat\varepsilon_i=y_i-\bm x_i\T\hat{\bm\beta}$.
(c) Cook's distance is $CD_i=\|\hat{\bm y}-\hat{\bm y}_{(i)}\|^2/(p\hat\sigma^2)$ with $\hat{\bm y}_{(i)}=X\hat{\bm\beta}_{(i)}$. Show that
$CD_i=\dfrac{\hat\varepsilon_i^2}{p\hat\sigma^2}\cdot\dfrac{H_{ii}}{(1-H_{ii})^2}$.
(d) Let $\hat\sigma^2_{(i)}$ be the unbiased estimate of $\sigma^2$ with case $i$ deleted. Show that
$\hat\sigma^2_{(i)}=\hat\sigma^2+\frac{1}{n-p-1}\bigl(\hat\sigma^2-\frac{\hat\varepsilon_i^2}{1-H_{ii}}\bigr)$.
\end{exercise}

\begin{exercise}[PS3 (2013), Problem 3(a)--(b): constrained least squares via QR]\label{ch05:ex-ps3-3ab}
Consider $H_0:\beta_{k+1}=\dots=\beta_p=0$. Let $X_I=[X_{[1]}\cdots X_{[k]}]$, $\hat{\bm\beta}_I=(X_I\T X_I)^{-1}X_I\T\bm y$ and
$\hat{\bm\beta}_0=(\hat{\bm\beta}_I,\bm 0_{p-k})$.
\begin{enumerate}[(a)]
  \item Show that $\hat{\bm\beta}_0$ minimises $(\bm y-X\bm\beta)\T(\bm y-X\bm\beta)$ subject to $\beta_j=0$ for $j>k$.
  \item With $X=QR$, show that $X_I=Q_IR_I$ where $Q_I$ is the first $k$ columns of $Q$ and $R_I$ the upper-left $k\times k$ block of $R$;
      verify $\hat{\bm\beta}_I=R_I^{-1}Q_I\T\bm y$ and $\hat{\bm y}_I=H_I\bm y$ with $H_I=Q_IQ_I\T$.
\end{enumerate}
\end{exercise}

\begin{exercise}[PS3 (2013), Problem 3(c)--(d): sequential ANOVA]\label{ch05:ex-ps3-3cd}
Let $A=\bigl(\begin{smallmatrix}Q\T\\ W\T\end{smallmatrix}\bigr)$ be orthogonal as in the proof of \cref{ch05:thm-normal} and $\bm z=A\bm y=(z_1,\dots,z_n)\T$.
(c) Prove that $\bm y\T\bm y=\sum_{i=1}^nz_i^2$, $\hat{\bm y}\T\hat{\bm y}=\sum_{i=1}^pz_i^2$, $\hat{\bm\varepsilon}\T\hat{\bm\varepsilon}=\sum_{i=p+1}^nz_i^2$, and for the constrained model
$\hat{\bm y}_I\T\hat{\bm y}_I=\sum_{i=1}^kz_i^2$, $\hat{\bm\varepsilon}_I\T\hat{\bm\varepsilon}_I=\sum_{i=k+1}^nz_i^2$.
(d) Using $\bm z_Q\sim N_p(R\bm\beta,\sigma^2I_p)$ independent of $\bm z_W\sim N_{n-p}(\bm 0,\sigma^2I_{n-p})$, deduce that
$SS_{\text{ERROR}}=\hat{\bm\varepsilon}\T\hat{\bm\varepsilon}\sim\sigma^2\chi^2_{n-p}$; that under $H_0$
$SS_{\text{REG}(k+1,\dots,p\mid1,\dots,k)}=\hat{\bm y}\T\hat{\bm y}-\hat{\bm y}_I\T\hat{\bm y}_I=\sum_{i=k+1}^pz_i^2\sim\sigma^2\chi^2_{p-k}$, independent of
$SS_{\text{ERROR}}$; and that $\hat F=\dfrac{SS_{\text{REG}}/(p-k)}{SS_{\text{ERROR}}/(n-p)}\sim F_{p-k,\,n-p}$. Arrange the results in the ANOVA
table with rows ``regression on $1,\dots,k$'' ($\hat{\bm y}_I\T\hat{\bm y}_I$, $k$ d.f.), ``regression on $k+1,\dots,p$ adjusting for
$1,\dots,k$'' ($p-k$ d.f., mean square $MS_0$, $\hat F=MS_0/MS_{\text{Error}}$), ``error'' ($n-p$ d.f.) and ``total'' ($\bm y\T\bm y$,
$n$ d.f.).
\end{exercise}

\printsolutions
